\documentclass{article}
\usepackage[utf8]{inputenc}
\usepackage[T1]{fontenc}
\usepackage{amsmath,amssymb,amsthm,bbm,float,graphicx,geometry,lmodern,mathtools,parskip,setspace,subcaption,diagbox}
\usepackage[colorlinks]{hyperref}
\usepackage{mathrsfs}
\usepackage{enumerate}
\usepackage{nicefrac}
\usepackage{todonotes}
\usepackage{comment}
\usepackage[notref,notcite]{showkeys}
\newcommand{\e}{\mathrm e}

%%%Tikz
\usepackage{tikz}
\usetikzlibrary{backgrounds}
\usetikzlibrary{patterns}
\usetikzlibrary{positioning, shapes.geometric}

%%% Captionsize
\usepackage{caption}
\captionsetup{font=footnotesize}
\usepackage{graphicx}
%%%%%%Formatting%%%%%%
\graphicspath{{Images/}}%put any images in this file
\parindent0pt%paragraph indent
\allowdisplaybreaks%breaks long equations onto multiple pages is needed
%%%%%%Commands%%%%%%
\DeclarePairedDelimiter{\set}{\lbrace}{\rbrace}
\newcommand{\N}{\mathbb{N}}
\DeclareMathOperator{\argmax}{arg \ max}
%%%%%%Environments%%%%%%

\newtheorem{theorem}{Theorem}[section]
\newtheorem{lemma}[theorem]{Lemma}
\newtheorem{corollary}[theorem]{Corollary}
\newtheorem{prop}[theorem]{Proposition}
\newtheorem{conj}[theorem]{Conjecture}

\theoremstyle{definition}
\newtheorem{definition}[theorem]{Definition}
\newtheorem{remark}[theorem]{Remark}

%%%%%%%% symbols and short cuts%%%%%%%
\renewcommand{\P}{\mathbb{P}}
\newcommand{\x}{\mathbf{x}}
\newcommand{\y}{\mathbf{y}}
\newcommand{\1}{\mathbbm{1}}
\newcommand{\G}{\mathscr{G}}
%\renewcommand{\H}{\mathscr{H}}
\newcommand{\C}{\mathscr{C}}
\newcommand{\X}{\mathbf{X}}
\newcommand{\E}{\mathbb{E}}
\newcommand{\0}{\mathbf{o}}
\newcommand{\R}{\mathbb{R}}
\renewcommand{\d}{\mathrm{d}}
\newcommand{\z}{\mathbf{z}}
\newcommand{\w}{\mathbf{w}}
\renewcommand{\v}{\mathbf{v}}
\renewcommand{\u}{\mathbf{u}}
\newcommand{\Z}{\mathbb{Z}}
\newcommand{\T}{\mathcal{T}}
\newcommand{\U}{\mathcal{U}}
\newcommand{\cX}{\mathcal{X}}
\newcommand{\cG}{\mathcal{G}}


%Commands Bas
\newcommand*{\be}{\begin{equation}}
\newcommand*{\ee}{\end{equation}}
\newcommand*{\ba}{\begin{aligned}}
\newcommand*{\ea}{\end{aligned}}

\newcommand{\D}{\mathbb D}
\newcommand{\cB}{\mathcal B}
\newcommand{\cP}{\mathcal P}
\newcommand{\cU}{\mathcal U}
%\newcommand{\ind}{\mathbbm 1}
\newcommand{\wt}{\widetilde}
\newcommand{\cR}{\mathcal R}
\newcommand{\sm}{\mathrm{sm}}
\newcommand{\basc}[1]{{\color{blue}{ \bf [~Bas:\ }\emph{#1}\textbf{~]}}}

%Packages Bas
%\usepackage{enumitem}
\usepackage{bbm}

%%%%%%Bibliography%%%%%%
\usepackage[style = numeric, sorting=nyt, url = false, abbreviate=false, maxbibnames=9, sortcites=true, doi = true, backend = biber, giveninits = true]{biblatex}
\renewbibmacro{in:}{\ifentrytype{article}{}{\printtext{\bibstring{in}\intitlepunct}}}
\bibliography{bib.bib}

\title{Percolation in lattice $k$-neighbor graphs}
\author{Benedikt Jahnel, Jonas K\"oppl, Bas Lodewijks, Andr\'as T\'obi\'as}
\date{\today}

\begin{document}

\maketitle

\begin{abstract}TTT
\end{abstract}

\section{Introduction}
The study of $k$-nearest neighbor type models in continuum percolation was initiated by Häggström and Meester \cite{HaggMee96}, who investigated the so-called \emph{undirected} $k$-nearest neighbor graph. Here, the vertex set of the random graph under consideration is a homogeneous Poisson point process in $\R^d$, and one connects two points by an edge if at least one of them is among the $k$ nearest neighbors of the other. One calls a maximal connected component in this graph a \emph{cluster}, and one says that the graph \emph{percolates} if it exhibits an infinite (equivalently, unbounded) cluster. Note that the probability of percolation is in $\{0,1\}$ thanks to ergodicity, and whether the graph percolates does not depend on the spatial intensity of the Poisson point process. Hence, the only parameters of the model are $k$ and $d$. It was showed in~\cite{HaggMee96} that this model does not percolate for any dimension if $k = 1$, but it percolates for $k=2$ if $d$ is large enough. Moreover, they proved that for any $d \geq 2$, there exists $k$ such that the graph percolates.

Balister and Bollobás \cite{BalisterBollobas13} investigated three possible senses of percolation in the model where each vertex is connected  to its $k$ nearest neighbors by a \emph{directed} edge: \emph{in-percolation} (resp.\ \emph{out-percolation}), which occurs by definition whenever some point of the point process exhibits an infinite incoming (resp.\ outgoing) path ending (resp.\ starting) at it, and \emph{strong percolation}, which means that there exists a strongly connected component in the graph, i.e., a component where from any point there exists a directed path to any other point. Additionally, they introduced another undirected graph, called the \emph{bidirectional} $k$-nearest neighbor graph, where one connects two vertices whenever they are mutually among the $k$ nearest neighbors of each other. For the two-dimensional case, they verified percolation in the undirected graph for $k \geq 11$, in the directed graph in all the three senses (in-, out- and strong percolation) for $k \geq 13$ and in the bidirectional graph for $k \geq 15$. It is easy to see that these graphs do not percolate in one dimension for any $k$. We observe that it follows from the proof of \cite[Theorem 2.2]{HaggMee96} that for any $d\geq 2$, for $k \geq 2$ sufficiently large, also the bidirectional graph percolates (despite the fact that the notion of bidirectional $k$-nearest neighbor graph does not appear in that paper).

%These $k$-nearest neighbor percolation models can be seen as an alternative of modeling connectivity in a telecommunication network, as opposed to the classical Gilbert graph, where one connects two points of the Poisson point process by an edge whenever their Euclidean distance is less than a given radius $r>0$. The latter graph was introduced by Gilbert (1961), who showed that for $d =2$, there is a phase transition in the spatial intensity $\lambda$ of the Poisson point process: There exists a critical intensity $\lambda_{\rm c}>0$ such that for $\lambda<\lambda_{\rm c}$, the graph almost surely consists of finite components, while for $\lambda>\lambda_{\rm c}$, the graph \emph{percolates}, i.e., it contains an infinite (equivalently, unbounded) connected component. While the Gilbert graph has unbounded degrees but bounded edge lengths, $k$-nearest neighbor type models have bounded degrees but arbitrarily long edges. A more realistic model for telecommunications, which is a proper subgraph of both the bidirectional $k$-nearest neighbor graph for some $k$ and the Gilbert graph for some $r$ is the so-called signal-to-interference plus noise ratio (SINR) graph, see \color{red} references \color{black}.

\cite{BalisterBollobas13} showed that for $d=2$, with high confidence, the undirected $k$-nearest neighbor graph percolates if and only if $k \geq 3$, in-, out- and strong percolation in the directed graph occurs if and only if $k\geq 4$, and the bidirectional graph percolates if and only if $k\geq 5$. By ``with high confidence'' they mean the following: They show that if a certain deterministic integral exceeds a certain deterministic value, then their aforementioned assertions are true. However, they were only able to evaluate these integrals via Monte Carlo methods, which yield an extremely small but positive error probability. Recently, \cite{JT22} showed that in the bidirectional graph there is no percolation for $k=2$ in any dimension, even if the underlying point process is not a Poisson point process but an arbitrary deletion-tolerant point process (satisfying some basic nondegeneracy conditions). Their proof exploits the simplicity of the structure of the bidirectional graph for $k=2$, which has degrees bounded by 2. Proving absence of percolation for $k=3,4$ seems to be out of reach at the moment even in the Poisson case, and it is also not entirely clear whether these assertions hold in all dimensions or only for $d=2$.

To understand the structure and elementary properties of $k$-nearest neighbor graphs better, in the current paper we take a step back and introduce a discrete counterpart of the (undirected, directed or bidirectional) $k$-nearest neighbor model. Here, $k$ is at most $2d$ (where $d$ still denotes the dimension), and the directed  $k$-nearest neighbor graph ($k$-DnG) is defined via connecting each vertex of $\Z^d$ to precisely $k$ of its $2d$ nearest neighbors (w.r.t.\ $\ell^1$-distance), which are chosen independently and uniformly for each vertex. If we now put an undirected edge between all nearest neighbor pairs connected by at least one edge (resp.\ by two directed edges) of the $k$-DnG, we obtain the undirected (resp.\ bidirectional) $k$-nearest neighbor graph, abbreviated as $k$-UnG (resp.\ $k$-BnG). In the $k$-DnG we study the analogue of out-percolation, which is often the most convenient choice, but in- or strong percolation could also be investigated similarly. Percolation in the $k$-DnG in our sense also implies percolation in the $k$-UnG, while percolation in the $k$-BnG graph implies percolation in the $k$-DnG. 

The motivation for studying this particular model is as follows. Consider the continuum nearest neighbor percolation model with $k=2$, where the intensity of the underlying Poisson point process such that the expected number points in a unit ball equals one, and let $Y_i$ denote the position of the $i^{\text{th}}$ nearest neighbor for $i\in\{1,2\}$. H\"aggström and Meester show that $|Y_i|$ converges in probability to one as $d$ tends to infinite, as well as that the conditional distribution of $Y_i$, given $|Y_i|=r$, is uniform on the sphere $\{x\in \R^d: |x|=r\}$~\cite[Lemma $3.2$]{HaggMee96}. Moreover, the volume of the intersection of two spheres with radius $r_1,r_2\in(0.9,1.1)$, whose centers are at least $0.9$ units apart, is negligible compared to the volume of either sphere~\cite[Lemma $3.3$]{HaggMee96}. As a result, for large dimensions the continuum nearest neighbor graph has connections at distance around $1$, which are established almost independently among different pairs of vertices. As such, an approximation of this graph on $\Z^d$ should yield similar behavior when $d$ is large. 

Note also that the $k$-BnG is a 1-independent Bernoulli bond percolation model where each edge is ``open'' (i.e., included in the $k$-BnG) with probability $p=k^2/(4d^2)$. That is, any fixed edge is open with probability $p$, and the events that some edges are open are independent whenever the edges are pairwise non-adjacent. The same holds for the $k$-UnG with $p=(1-k/(4d))k/d$. Let us also mention that these lattice $k$-nearest-neighbor graphs present basic examples of {\em negatively correlated} random fields, in contrast to many of the classical models like the random cluster model, see \color{red}Grimmett 2006\color{black}. Indeed, note that the presence of an edge makes the presence of neighboring edges more unlikely. 

Our main results include results for the presence and absence of percolation for all the three variants of the model, and their proofs use different ideas and/or prior results. Again, the case $d=1$ is trivial: There is no percolation in any of the three senses if $k=1$, and all (directed resp.\ undirected) nearest-neighbor edges of $\Z$ are in the graph if $k=2$. Also abserve that there is a natural ordering of the models with respect to percolation in the sense that existence of percolation in the bidirectional model implies existence of percolation in the directed model, which implies percolation in the undirected mode. The following results are meant for $d \geq 2$; for a summary of the ones among them corresponding to percolation in the directed and undirected graphs, see also Table~\ref{table-UD}. 
\begin{itemize}
    \item For $k=1$, our Proposition~\ref{prop-U1} tells that even the undirected graph does not percolate for any dimension. The proof of this assertion involves a mass-transport type argument, which is similar in spirit to the analogous proof in the continuum case by \cite{HaggMee96}, combined with a classical estimate for the number of self-avoiding paths (following \cite{Grimmett1999}). Let us note that proving just that the $1$-DnG does not percolate is simpler (see the proof of Lemma~\ref{lem:surenoperc}).
    \item For $k=2$, we show in Theorem~\ref{theorem-U22} that the $k$-UnG percolates for $d=2$. The proof of this assertion is an adaptation of the classical Peierls argument involving a dual lattice in the spirit of \color{red} Durrett 88 and \cite{JanSwart} \color{black}. The percolation for $k=2$ in the $k$-DnG for $d \geq 2$ and in the $k$-UnG for $d \geq 2$ remains open, and we will comment on this later during this introduction.
    \item For $k \geq d+1$, our Lemma~\ref{lem:sureperc} shows that percolation occurs in a very strong sense: There is an infinite outgoing path starting from any fixed vertex with probability 1. The proof of this assertion is elementary and not based on any prior results.
    \item For $k=3$ we have percolation if $d \geq 5$ and for $k \geq 4$ if $d \geq 2$ (here, the cases $d=2,3$ follow thanks to the previous bullet point). This is the content of Theorem~\ref{thm-CoxDurrett}, whose proof is based on arguments by \cite{CoxDurrett} on oriented percolation in at least 4 dimensions, combined with some explicit computations in a few low-dimensional cases where the estimates by \cite{CoxDurrett} are not strong enough yet.
    \item Using the recent improved lower bounds by \cite{improvedLB} on the smallest $p$ such that any 1-independent Bernoulli bond percolation model percolates, in Proposition~\ref{prop-1indep} we obtain percolation in the $k$-BnG for $k>2\sqrt{0.8457}d \approx 1.8392d$, and for $d$ sufficiently large even for $k > 2\sqrt{0.5847}d \approx 1.5293d$. In particular, the bidirectional graph with $k=2d-1$ percolates for $d \geq 7$ and the one with $k=2d-2$ for $d \geq 13$. On the other hand, using negative correlations of the $k$-BnG, we show that if $k(k-1) < 2d (2d-1)/c(d)$, where $c(d)$ is the connective constant of $\Z^d$ (whose definition we will recall in the beginning of Section~\ref{sec:B}), then there is no percolation in the bidirectional graph. This implies that if the $k$-BnG percolates, then $k=\Omega(\sqrt d)$ as $d\to\infty$. In particular, for fixed $k$ there is no percolation if $d$ is sufficiently large. Such a phenomenon has not been observed in the continuum setting so far. It also follows that for $k=2$ we have no percolation in the $k$-BnG for any $d\geq 2$, which is the analogue of the aforementioned assertion by \cite{JT22} in the continuum case.
\end{itemize}

\begin{table}
\centering
\begin{tabular}{|l|c|c|c|c|c|c} 
 \hline
\diagbox{$d$}{$k$} & 1 & 2 & 3 & 4 & $\geq 5$ \\ \hline 
1 & no & yes & - & - & - \\ \hline
2 & \color{blue}no\color{blue} & \color{pink}U yes\color{black}, \color{gray}D unknown & \color{green}yes\color{black} & yes & - \\ \hline
3 & \color{blue}no\color{blue} & \color{gray}unknown & \color{gray}unknown & \color{green}yes\color{black} & \color{green}yes\color{black} \\ \hline
4 & \color{blue}no\color{blue} & \color{gray}unknown & \color{gray}unknown & \color{red}yes\color{black} & \color{green}y\color{red}es\color{black} \\ \hline
$\geq 5$ & \color{blue}no\color{blue} & \color{gray}unknown & \color{red}yes & \color{red}yes & \color{red}yes \\  \hline
\end{tabular}

\caption{Percolation results and open cases for the $k$-UnG and the $k$-DnG. In the table, for given $k,d$, `yes' means that both graphs percolate and `no' that none of them does, while `unknown' (without anything else in the cell) means that both cases are open. The coloring subdivides the cases according to the assertions from which they follow. Black: trivial, grey: open, blue: Proposition~\ref{prop-U1} (for the directed case, Lemma~\ref{lem:surenoperc} also suffices)\color{black}, pink: Theorem~\ref{theorem-U22}, green: Lemma~\ref{lem:sureperc},  red: Theorem~\ref{thm-CoxDurrett} (and green-and-red: both of the latter two).}\label{table-UD}
\end{table}

We see that some of our main results are similar to known assertions and unproven conjectures for the continuum $k$-nearest neighbor graphs, which indicates that despite the lack of the unbounded-range dependencies that are characteristic for the continuum models, our models are indeed their natural discrete analogue. Even without being motivated by the continuum case, our graphs show interesting phase transitions and have many connections to the existing literature (as listed above), while their definition is very simple, which gives a reason to study them on their own right. Nevertheless, to our best knowledge, our graphs have not been introduced before, which we found somewhat surprising. Similar discrete models such as corner percolation \cite{cornerperc} and constrained-degree percolation \cite{constraineddegreeperc} are known from the literature, but they all seem to be significantly different from ours \color{red} want to say more? \color{black} and they have not proven to be useful to verify (absence of) percolation in our models so far.

The present paper leaves some cases open, which would be interesting to resolve in the future. As already mentioned, we do not know whether the 2-DnG percolates in any dimension, but numerical simulations suggest that this actually seems to be the case for $d=2$ \color{red} want to add simulations? \color{black}. Such an assertion would be particularly interesting because the nearest-neighbor edges of $\Z^d$ missing from the $d$-DnG also form a $d$-DnG, i.e., if the $d$-DnG percolates, we have a coexistence of infinite clusters of chosen and non-chosen edges, which is generally considered as untypical for a two-dimensional percolation model. While we expect that $k=2$ is critical in all dimensions higher than one, we also have two ``sporadic'' cases where the question of percolation in the directed graph is still open: $k=d=3$ and $k=3,d=4$. We prove that if one can show absence of percolation for $k=d=3$, then the same assertion follows for $k=d=2$. As it is not clear whether percolation for some $k$ and $d$ implies percolation for the same $k$ and for all $d' \geq d$ (for the same type of graph), $k=2$ is also open in the undirected case for $d \geq 3$. For the $k$-BnG, we have seen that for any $k$ we have absence of percolation if $d$ is large enough, but finding the precise order of the smallest $k$ for which we have percolation as $d \to \infty$ is not known yet; according to our results it could be between constant times $\sqrt d$ and constant times $d$. Moreover, we do not have any positive results for the $k$-BnG if the dimension is less than 7. One would (at least naively) imagine that choosing $k=2d-1$ nearest neighbors in the bidirectional sense should be sufficient for percolation even in some lower dimensions, especially because the nearest-neighbor edges of $\Z^d$ missing from the $2d-1$-BnG form a $1$-UnG (cf.~Section~\ref{sec:phases} for more details), about which we have showed that it does not percolate.

The remainder of this paper is organized as follows. In Section~\ref{sec:set} we provide a formal definition of our $k$-DnG, $k$-BnG and $k$-UnG models, and we present our results corresponding to these graphs in Sections~\ref{sec:D}, \ref{sec:B} and \ref{sec:U} along with their proofs, respectively. In Section~\ref{sec:discussion} we discuss our results and mention some related conjectures and open questions. The proofs of some technical results \color{red}(really plural?) \color{black} are carried out in the Appendix, we will explained where to find the proofs after presenting the results.

\section{Setting and main results}\label{sec:set}

Consider the $d$-dimensional cubic lattice $\mathbb Z^d$. We are interested in the percolation behavior of the {\em $k$-neighbor graph} ($k$-nG) in which each vertex independently chooses uniformly $k\le 2d$ of its $2d$ nearest neighbors. We distinguish three different types of $k$-nGs:
\begin{enumerate}
\item The {\em directed $k$-nG} ($k$-DnG) in which we open a directed edge from the vertex to each of the $k$ chosen neighbors. 
\item The {\em undirected $k$-nG} ($k$-UnG) in which we open an undirected edge from the vertex to each of the $k$ chosen neighbors.
\item The {\em bidirectional $k$-nG} ($k$-BnG) in which we open an undirected edge between two vertices whenever both choose the other one as one of its $k$ chosen neighbors.
%\item The {\em exclusively unidirectional $k$-nG} or ($k$-XnG) in which we open an undirected edge between two vertices whenever exactly one of them chooses the other one as one of its $k$ chosen neighbors. The letter X refers to the ``XOR'' (exlusive ``or'') in the edge-drawing rule.
\end{enumerate}

We are interested in the behavior of the percolation probabilities 
$$
\theta^{\rm D}(k,d)=\mathbb P(o\rightsquigarrow\infty \text{ in the $k$-DnG on $\mathbb Z^d$}),
$$
respectively $\theta^{\rm U}(k,d)$ and $\theta^{\rm B}(k,d)$, %and $\theta^{\rm X}(k,d)$, 
where $o\rightsquigarrow\infty$ represents the event that there exists a path along open edges from the origin to infinity. Let us note that the percolation probabilities are clearly non-decreasing functions in $k$. 

Let us mention that the outdegree in $k$-DnG is $k$ almost surely. However, the degree in $k$-BnG is a binomial random variable with parameters $k$ and $k/2d$, i.e., the expected degree is given by $k^2/2d$ which is less than $k$, unless $k=2d$. In case of the $k$-UnG, the degree is distributed according to $k+D$, where $D$ is a binomial random variable with parameters $2d-k$ and $k/2d$, and hence, the expected degree is given by $k(4d-k)/(2d)$.

\subsection{Results for the directed $k$-neighbor graph}\label{sec:D}
\begin{lemma}\label{lem:surenoperc}
We have for all $d\ge 1$ that $\theta^{\rm D}(1,d)=0$.
\end{lemma}
\begin{proof}
Indeed, if there exists a path starting from the origin and reaching an $\ell_1$-distance $n$ from the origin, then this path is unique and has at least $n$ steps. But such a path exists with probability at most $(2d)^{-n}$, which converges to zero as $n$ tends to infinity.
\end{proof}
%Note that the $k$-nnG has a natural dual in which each existing directed edge is removed and all missing edges are present. This graph is edge-disjoint but not vertex disjoint and it is in distribution equivalent to the $(2d-k)$-nnG. 
% \begin{lemma}\label{lemma2}
% We have that $\theta(2d-1,d)>0$.
% \end{lemma}
% \begin{proof}
% Note that by ergodicity, since $\theta(1,d)=0$ in the $(2d-1)$-nnG in $d$ dimensions, there exist only finite components in the $1$-nnG in $d$ dimensions. But this means that the $(2d-1)$-nnG in $d$ dimensions percolates and thus $\theta(2d-1,d)>0$. 
% \end{proof}
\begin{lemma}\label{lem:sureperc}
    We have that $\theta^{\rm D}(k,d)=1$ for all $k\ge d+1$.
\end{lemma}
In particular, from Lemma~\ref{lem:surenoperc} and Lemma~\ref{lem:sureperc} we see that for $d=2$ the critical $k$ is given by $2$. Simulations suggest that $\theta^{\rm D}(2,2)>0$. We consider a proof of this statement a major success. 
\begin{proof}
    We will use a growth argument and show that the maximal distance to the origin is strictly increasing between generations. To make this precise, denote by $G_n$ the new vertices discovered in the $n$-th step, where we start with $G_0 = \{0\}$ and at every step $n \to n+1$, each vertex in $G_n$ chooses $k$ of its neighbours uniformly at random as potential successors and $G_{n+1}$ is then the set of all potential successors that have not been discovered before for any $m  \leq n$. For $x\in\Z^d$, we let $\|x\|_1$ denote the $\ell_1$-distance between $x$ and the origin. We show by induction that 
    \begin{align}\label{eqn:growth-property}
         \forall n  \in \N:   \quad \max_{x \in G_n}\|x\|_1 < \max_{x \in G_{n+1}}\|x\|_1.  
    \end{align}
    This then implies that there exists an infinitely long directed path starting at the origin. 
    
    For $n=0$ the inequality in~\eqref{eqn:growth-property} is clear. For the induction step, let $x \in G_n$ be a vertex (possibly non unique) that achieves the maximum. Then $x$ has $d$ neighbours $y$ with $\|x\|_1 < \|y\|_1$. If $k  \geq 2d-d+1$, then we necessarily have to choose one of these $y$ as a potential vertex for the new generation. This $y$ cannot have been in any of the previous generations by the induction hypothesis. Therefore we have $y \in G_{n+1}$ and~\eqref{eqn:growth-property} follows. 
    \end{proof}

Let us turn our attention towards the supercritical percolation phase. We have the following result. For this, consider two independent random walks $S$ and $S'$ with $S_0=e_1$ and $S'_0=e_2$, where $(e_1,\dots, e_d)$ denote the unit vectors. The walks move as simple random walks in the positive direction, i.e., $S_{n+1}=S_n+e_i$ where $i$ is uniformly chosen from $\{1,\dots ,d\}$. We define by
\begin{align*}
\rho(d)=\P(\exists m\ge 1\colon S_m=S'_m)\in(0,1]
\end{align*}
the meeting probability of the walkers. 

\begin{theorem}\label{thm-CoxDurrett}
$\theta^{\rm D}(k,d)>0$ holds if $d \geq 2$ and $k \geq 4$, or if $d \geq 5$ and $k=3$. 
\end{theorem}
\begin{proof}
We follow \cite[Section 2]{CoxDurrett}. Let $\mathcal R_n$ denote the set of directed paths from the origin to level $n$ (that is, all vertices at $\ell_1$ distance $n$) with strictly increasing $\ell_1$ distance \color{red} OK? \color{black}, and let $N_n$ be the number of open paths in $\mathcal R_n$. Then $W_n=(2/k)^nN_n$ is a martingale with respect to $\mathcal F_n$, the sigma-algebra generated by the first $n$ steps of a simple random walk $T$ with $T_0=o$ and $T_{n+1}=T_n+e_i$ where $i$ is uniformly chosen from $\{1,\dots ,d\}$. Indeed, using the fact that the expected number of outgoing (i.e., w.r.t.\ $\ell_1$ distance increasing) edges  is given by the expectation of a hypergeometric random variable, i.e., 
\begin{align*}
\sum_{l=0}^k l\frac{\binom{d}{l}\binom{d}{k-l}}{\binom{2d}{k}}=\frac{k}{2},
\end{align*}

\begin{comment} 
\basc{Observe that this is only true for vertices $v$ such that the number of neighbours with larger (resp.\ smaller) $\ell_1$-distance equals $d$ (resp.\ $d$). There are, however, vertices for which this is not the case. In any dimension $d\geq 2$ and any $n\in\N$, there are always $2d$ vertices at distance $n$ such that they have $2d-1>d$ many neighbours at $\ell_1$-distance $n+1$ and only $1$ neighbour at $\ell_1$-distance $n-1$. In particular, when $d=k=2$,
\be 
N_1=k \text{ almost surely, and } \E[N_2]=12\Big(\frac{k}{2d}\Big)^2=3\neq 2. 
\ee
Similarly, 
\be \ba
\E[N_3]&=4\cdot 3\Big(\frac{k}{2d}\Big)^3+4\cdot 2\cdot 2\Big(\frac{k}{2d}\Big)^3=28\Big(\frac{k}{2d}\Big)^3=\frac72\neq 2,\\
\E[N_4]&=4\cdot 3\Big(\frac{k}{2d}\Big)^4 +8\cdot 3\cdot 2 \Big(\frac{k}{2d}\Big)^4=\frac{15}{4}\neq 2. {\color{red} \text{this may be incorrect}}
\ea \ee  
As $n$ grows, $\E[N_n]$ will become closer to $2$ (in the case $k=2$, or to $k(k/2)^n$ in general, see also my comment further on) I think, as the contribution of the paths from the $4$ vertices with $3>d$ many neighbours decreases relatively compared to the contribution of the other paths. \\
In general (with $d=2$ and $k\leq d$), at $\ell_1$-distance $n$ there are $4(n-1)$ vertices who have, in expectation, $k/2$ many neighbours that increase their $\ell_1$-distance w.r.t.\ their parent, and $4$ vertices with, in expectation, $3k/4$ many neighbours that increase their $\ell_1$-distance w.r.t.\ their parent.\\
As a result, $W_n$ is a sub-martingale rather than a martingale. 
}
\end{comment}
we can use independence to see that 
\begin{align*}
\E[N_{n+1}|\mathcal F_n]=\frac{k}{2}N_n.
\end{align*}
Since $\E[W_1]=1$, we can apply the martingale convergence theorem, to ensure the existence of a random variable $W$ with 
\begin{align*}
W_n\to W\qquad \text{almost surely}.
\end{align*}
Using the second-moment method it now suffices to show that $\limsup_{n\uparrow\infty}\E[W_n^2]<\infty$.

For this, note that 
\begin{align*}
\E[N_n^2]=\sum_{s,t\in \mathcal R_n}\P(s,t\text{ open})=\sum_{s,t\in \mathcal R_n}p^{K(s,t)}q^{L(s,t)}p^{2(n-K(s,t)-L(s,t))}
\end{align*}
where $p=k/2d$ denotes the probability that a given edge is open, $q=k(k-1)/(2d(2d-1))$ denotes the probability that two different edges that emerge from the same vertex are open, $K(s,t)$ is the number joint edges in two paths $s$ and $t$, and $L(s,t)$ is the number of vertices $x$ in $s$ and $t$ such that $s$ and $t$ do not have a joint edges after $x$.  Now note that $q<p^2$, which implies that
\[ \E[N_n^2] \leq \sum_{s,t\in \mathcal R_n} p^{2n-K(s,t)}. \]
The right-hand side is precisely the same as the right-hand side of the display below \cite[Equation~(2$\cdot$7)]{CoxDurrett}. Thus, a verbatim application of the arguments of \cite[Section 2]{CoxDurrett} implies that 
\[ \lim_{n\to\infty} \E W_n^2 <\infty \]
holds if and only if 
\[ p = \frac{k}{2d} < \varrho(d), \]
where for two independent (simple, symmetric, nearest neighbor) random walks $\widetilde S=(\widetilde S_n)_{n\in\N_0}$ and $\widetilde S'=(\widetilde S'_n)_{n\in\N_0}$ on $\Z^d$ started from the origin, we define
\[ \varrho(d) = \P(\exists m\geq 0 \colon \widetilde S_m = \widetilde S'_m \text{ and } \widetilde S_{m+1} = \widetilde S'_{m+1}). \]
Let further $\tau_d=\inf \{ m \geq 0 \colon \widetilde S'_m \text{ and } \widetilde S_{m+1} = \widetilde S'_{m+1} \}$.
According to \cite[p.~155]{CoxDurrett} we have for all $d$ that
\begin{equation}\label{firstfew} \P(\tau_d=0)=\frac{1}{d}, \qquad \P(\tau_d=1)=0, \qquad \P(\tau_d=2)=d^{-3}-d^{-4}, \end{equation}
and for $3 \leq k \leq d$,
\[ \P(\tau=k) \leq d^{-k} k!, \]
while for $j\geq 1$ and $k >jd$,
\begin{equation}\label{bigk} \P(\tau=k) \leq d^{-1} (2\pi d)^{1/2} \Big( \frac{\e^{-1/13}}{\sqrt{2\pi}} \Big)^d j^{-\frac{1-d}{2}}, \end{equation}
so that
\begin{equation}\label{CDUB} \begin{aligned} \varrho(d)=\P(\tau_d<\infty) &\leq \frac{1}{d} + \frac{1}{d^3} -\frac{1}{d^4} + \sum_{k=3}^d d^{-k} k! +  (2\pi d)^{1/2} \Big( \frac{\e^{-1/13}}{\sqrt{2\pi}} \Big)^d \sum_{j=1}^{\infty} j^{-\frac{1-d}{2}} \\ & = \frac{1}{d} + \frac{1}{d^3} -\frac{1}{d^4} + \sum_{k=3}^d d^{-k} k! +  (2\pi d)^{1/2} \Big( \frac{\e^{-1/13}}{\sqrt{2\pi}} \Big)^d \zeta\Big( \frac{d-1}{2} \Big), \end{aligned} \end{equation}
where $\zeta$ denotes the Riemann zeta function. (Note that since $\varrho(d)>1/d$, this approach cannot give any positive percolation result for $k=2$.)

%Therefore, it suffices to find the smallest $d$ such that the upper bound of~\eqref{CDUB} is less than $3/(2d)$, and this will imply that $\theta^{\rm D}(k',d') >0$ for all $k' \geq k$ and $d' \geq 3$.
%Note that this will not give any positive percolation result for $k=2$ because $k/2d = 1/d$ is smaller than the right-hand side of~\eqref{CDUB}. The next table shows the values of the right-hand side for $d=4,5,6,7$ and the respective smallest values of $k$ such that these bounds imply $\theta^{\rm D}(k,d)>0$.

The next table shows the smallest values of $k$ such that the right-hand side of~\eqref{CDUB} (which we computed numerically) is less than $k/(2d)$, for $d=4,5,6,7$.
\begin{table}
\begin{center}
    \begin{tabular}{|c|c|c|c|c|} \hline
       &  $d=4$ & $d=5$ & $d=6$ & $d=7$ \\ \hline
        Value of the r.h.s.\ of~\eqref{CDUB}   & 0.693093  & 0.394622  & 0.268615 & 0.199707  \\ \hline
         Smallest $k$ such that r.h.s.\ of~\eqref{CDUB} is less than $\tfrac{k}{2d} $ & 6 & 4 & 4 & 3 \\ \hline 
    \end{tabular}
\end{center}
\caption{Values of the right-hand side of~\eqref{CDUB} and the smallest $k$ such that this right-hand side is below $\tfrac{k}{2d} $, for $4 \leq d \leq 7$.}\label{table-CD}
\end{table}
%Via additional numerical computations we see that the right-hand side of~\eqref{CDUB} minus $k/(2d)$ is negative for any $k\geq 3$ if $7 \leq d \leq 10$ and for any $k\geq 4$ if $5 \leq d \leq 10$. 
Now, it is elementary to check that the right-hand side of~\eqref{CDUB} minus $k/(2d)$ stays negative for any $k \geq 3$ if $d \geq 7$, thanks to the following lemma:
\begin{lemma}\label{lemma-largedmon}
For any $d \geq 7$ and $k \geq 3$, we have
\begin{equation}\label{largedmon} \frac{1}{d} + \frac{1}{d^3} -\frac{1}{d^4} + \sum_{k=3}^d d^{-k} k! +  (2\pi d)^{1/2} \Big( \frac{\e^{-1/13}}{\sqrt{2\pi}} \Big)^d \zeta\Big( \frac{d-1}{2} \Big) - \frac{k}{2d} \leq 0.  \end{equation}
\end{lemma}
The proof of this lemma is carried out in Section~\ref{sec:largedmon} in the Appendix. Given that according to Table~\ref{table-CD}, \eqref{largedmon} also holds for $k=3$ and $d=7$, $k=4$ and $5 \leq d \leq 7$, as well as for $k=6$ and $d=4$, Lemma~\ref{lemma-largedmon} has the following consequences.  % \color{red} OK? For this it suffices to show that the rhs of~\eqref{CDUB} for dimension $d+1$ is at most $d/(d+1)$ times the same for dimension $d$, which I think is possible if $d \geq 10$. We just have to compensate the additional term $(d+1)^{-(d+1)} (d+1)!$ in the sum in dimension $d+1$. \color{black} 
First, $\theta^{\rm D}(5,d) >0$ for all $d \geq 5$. Moreover, Lemma~\ref{lem:sureperc} implies that $\theta^{\rm D}(5,d)=1$ holds for $d=3,4$. Thus, we have $\theta^{\rm D}(k,d) >0$ whenever $k\geq 5$ (given that it is well-defined, i.e., $k\leq 2d$). The table also tells that  $\theta^{\rm D}(4,d)>0$ for all $d \geq 5$. Moreover, $\theta^{\rm D}(4,2)=1$ is clear and $\theta^{\rm D}(4,3)=1$ follows from Lemma~\ref{lem:sureperc}. Finally, $\theta^{\rm D}(3,d)>0$ for $d \geq 7$ can be seen in the table (and $\theta^{\rm D}(3,2)=1$ is another consequence of Lemma~\ref{lem:sureperc}). Therefore, in order to finish the proof of the theorem, we have to treat three special cases separately: $(k,d)=(4,4)$, $(k,d)=(3,5)$ and $(k,d)=(3,6)$. 

As for the first one, we compute (using a computer) the precise values of $\P(\tau_4=k)$ for $k=3,4,5$ via generating all length-6 random walk trajectories and computing how many of them belong to these events respectively (for $\P(\tau_4=l)$, $l=0,1,2$, we of course obtain the same precise values as \cite{CoxDurrett}). This yields
\[ \P(\tau_4=3)=\frac{3}{512}, \qquad \P(\tau_4=4)=\frac{279}{65\ 536}, \qquad \P(\tau_4=5)=\frac{831}{262\ 144}. \]
Thus, using~\eqref{firstfew} and~\eqref{bigk},
\[ \begin{aligned} \P(\tau_4<\infty) & \leq \frac{1}{4} + \frac{3}{256} + \frac{3}{512} + \frac{279}{65\ 536} + \frac{831}{262\ 144} + \frac{3}{4} \sqrt{8\pi} \Big( \frac{\e^{-1/13}}{\sqrt{2\pi}} \Big)^4 + \sum_{j=2}^{\infty} \sqrt{8\pi} \Big( \frac{\e^{-1/13}}{\sqrt{2\pi}} \Big)^4 \sum_{j=1}^{\infty} j^{-\frac{3}{2}} \\ & \approx 0.495542 < 0.5 = \frac{k}{2d}. \end{aligned}\]
It follows that $\theta^{\rm D}(4,4)>0$.

With similar methods we can show that $\theta^{\rm D}(3,5)>0$ and $\theta^{\rm D}(3,6)>0$. Indeed, for $d=5$ we obtain via explicit computer computation
\[ \P(\tau_5=3)=\frac{44}{15\ 625}, \qquad \P(\tau_5=4)=\frac{712}{390\  625}, \qquad \P(\tau_5=5)=\frac{12\ 136}{9\ 765\ 625}. \]
This already implies that  $\theta^{\rm D}(3,5)>0$ because
\[ \P(\tau_5<\infty) \leq \frac{1}{5} + \frac{4}{625} + \frac{44}{15\ 625} + \frac{712}{390\  625} + \frac{12\ 136}{9\ 765\ 625} + \sqrt{10\pi} \Big( \frac{\e^{-1/13}}{\sqrt{2\pi}} \Big)^5 \zeta(2) \approx 0.275703, \]
%\color{red} (for fun we could write that the last summand is $\sqrt 5/(24 \e^{5/13})$) \color{black} 
which is less than $k/(2d)=0.3$. For $d=6$, determining the value of $\P(\tau_6=k)$, $k=3,\ldots,6=d$ numerically would be computationally quite demanding. However, we see from the table that the right-hand side of~\eqref{CDUB} is just slightly larger than $k/(2d)=1/4=0.25$. It turns out that in this case, it suffices to compute $\P(\tau_6=3)$ numerically, which equals $\frac{35}{23328}$. Estimating $\P(\tau_6>3)$ according to~\eqref{CDUB}, we obtain
\[ \P(\tau_6<\infty) \leq \frac{1}{6} + \frac{5}{1296} + \frac{35}{23\ 328} + \sum_{k=4}^6 6^{-k} k! + \sqrt{12\pi} \Big( \frac{\e^{-1/13}}{\sqrt{2\pi}} \Big)^6 \zeta(5/2) \approx 0.242338 < 0.25=\frac{k}{2d}. \]
\end{proof}
\color{red} What follows until Theorem~\ref{theorem-3322} is to be shortened/removed/updated. The conjecture after the proof of Theorem~\ref{theorem-3322} too. \color{black}

\color{green} After this, the remaining unresolved cases are: $k=2$ and $d \geq 2$; $k=3$ and $d=3,4$. For $(k,d)=(3,4)$ it is not impossible that some Cox--Durrett style argument works, but so far I only see a solution via computing $\P(\tau_4=k)$ until ca.\ $k=12$, which is too much for our numerical capacity. For $d=3$ and for $k=2$ the Cox--Durrett (or Bene's similar) argument doesn't work. \color{black}


\bigskip
It would be desirable to proof the following conjecture.
\begin{conj}
For all $k\ge 1$ we have that, whenever $\theta^{\rm D}(k,d)>0$, then also $\theta^{\rm D}(k,d+1)>0$. 
\end{conj}
If this is true, then $k=2$ would always be critical. It is not to be expected that $d\mapsto\theta^{\rm D}(k,d)$ is non-decreasing since $\theta^{\rm D}(3,2)=1$ but heuristically $\theta^{\rm D}(3,3)<1$. IDEA: USE COUPLING. IS $n\mapsto \theta^{\rm D}(k+n,d+n)$ NON-DECREASING? 

\begin{theorem}\label{theorem-3322}
For all $k,d\ge1$ we have that $\theta^{\rm D}(k+1,d+1)\ge \theta^{\rm D}(k,d)$.
\end{theorem}
\begin{proof}
We use a randomized coupling approach. Let us write $x=(x_1,x_2,x_3)\in \Z^3$ and $x'=\pi(x)=(x_1,x_2)\in \Z^2$ for the projection of $x$ to its first two coordinates. By $\omega_{x}$ we denote the three-dimensional three-neighbor directed edge configuration of $x\in \Z^3$ and by $\omega'_{x'}=\pi(\omega_{x})$ the associated edge configuration of $x'$ in $\Z^2$, in words we forget all edges facing up or down and only keep all remaining edges which we consider as part the edge set of $\Z^2$. 

We further write $\bar x=(x',n)_{n\in\Z}$, for the vector of nodes {\em above and below} $x'\in \Z^2$ in the third dimension. The idea now is to define probability kernels $\nu_x(\cdot|\bar x)$, step-by-step with respect to $x$, in such a way that the joined \color{red} András: joint? \color{black} kernel has two properties:
\begin{enumerate}
\item Integrating the joined kernel with respect to $\omega$ gives the distribution of the two-dimensional directed 2NN graph. 
\item If the connected component of the origin in $\omega$ is finite, then the joined kernel only puts positive mass on finite components of the origin in two dimensions as well. 
\end{enumerate}

To make this precise, we need to consider ordered pairs of neighbors of nodes $x'\in \Z^2$ with some third coordinate, i.e., we write $D_{x'}=\{((a,n),(b,m))\colon n,m\in \Z, a\neq b\in \Z^2, \text{ with }a\sim x'\text{ and }b\sim x'\}$, where $\sim$ represents the nearest-neighbor relation, here in $\Z^2$. Then, we start with the origin $o\in \Z^2$ and define for $\bar o$ the probability kernel $\nu_{o,0}(\cdot|\bar o)$ as a probability measure on $D_o$ as follows. 
\begin{enumerate}
\item If $|\pi(\omega_{(o,0)})|\ge 2$, pick a first neighbor $a_o$ uniformly at random from the available connected neighbors in $\Z^2$ and the a second neighbor $b_o$ uniformly at random from the remaining available connected neighbors in $\Z^2$. Write for the outcome $((a_o,0)(b_o,0))$.
\item If $|\pi(\omega_{(o,0)})|=1$, 
keep the single available connected neighbor in $\Z^2$ as $a_o$. Then, randomly choose a direction up or down (there must exist a directed edge facing up as well as down) and follow these steps:
\begin{enumerate}
\item If $|\pi(\omega_{(o,\pm 1)})\setminus \pi(\omega_{(o,0)})|\ge 1$, 
pick the second neighbor $b_o$ uniformly at random from the remaining available connected neighbors in $\Z^2$. Write for the outcome $((a_o,0)(b_o,\pm 1))$.
\basc{Do we think of $\pm 1$ as the point $(0,0,\pm 1)$ here? If not, I'm not entirely sure what is meant with $(b_o,\pm 1)$.}
\item If $|\pi(\omega_{(o,\pm 1)})\setminus \pi(\omega_{(o,0)})|= 0$, follow the existing arrow in the same direction as in step 2.~and repeat from (a) with $\omega_{(o,\pm 1)}$ replaced by $\omega_{(o,\pm 2)}$.
\end{enumerate}
\end{enumerate}
Now, almost surely, the construction ends and, due to the complete symmetry in the construction, $\nu_o$ reproduces the uniform two-nearest-neighbor distribution for the origin. That is, for $E^2_x=\{\{a,b\}\colon a\neq b\in \Z^2\text{ with }a\sim x\text{ and }b\sim x\}$ and $F\colon D_{x'}\to E^2_x$, $((a,n),(b,m))\mapsto \{a,b\}$ and $U^i_x$ the uniform distribution on $E^i_x$, then, for all $\omega'_o\in E^2_o$,
\begin{align*}
\int \Big(\bigotimes_{n\in \Z}U^3_{(o,n)}\Big)(\d \bar \omega_{\bar o})\nu_{o,0}(F^{-1}(\omega'_o)|\omega_{\bar o})=U^2_o(\omega'_o).
\end{align*}

Let us describe the following steps in words. Under $\nu_{o}(\cdot|\omega_{\bar o})$ we are presented with two neighbors of the origin in $\Z^2$ that both also carry a third coordinate. The third coordinate represents the level at which the neighbor is discovered. Starting with $(a_o,n)$, the first discovered neighbor, we can sample its neighbors, using the same algorithm, based on the information provided by $\bar a_o$, now started at level $n$ (and not at level $0$). The same can be done for the second neighbor $(b_o,m)$. In this fashion, we can step-by-step explore the connected component of the origin in $\Z^2$, where in any step, whenever we discover a new vertex $y\in \Z^2$, we only use information from the associated vector $\bar y$, 

\begin{align*}
\nu(\d \omega'|\omega)=\int \nu_{(o,0)}(\d(\bar a,\bar b)|\bar \omega_{\bar o})\int\nu_{\bar a}(\d(\bar c_{\bar a},\bar d_{\bar a})|\bar \omega_{\bar a})\int\nu_{\bar b}(\d(\bar e_{\bar b},\bar f_{\bar b})|\bar \omega_{\bar b})\cdots \1\{F(\bar a \bar b \bar c \bar d \bar e \bar f\dots)=\omega'\}
\end{align*}

Now, the key point is the following. Imagine that the algorithm presents us a configuration of open edges in $\Z^2$ such that the connected component of the origin reaches the boundary of a box with side-length $n$. Then, necessarily, any configuration in $\Z^3$ that is used as an input for the algorithm must also have the property that the origin is connected to the boundary of a box of side-length $n$, now of course in  $\Z^3$. Hence, denoting by $A^i_n$ the event that the origin is connected to the boundary of the centered box of side-length $n$ in $\Z^i$  \color{red} András: maybe we should also say that the boxes should be centred around the origin \color{black} and by $\P_i$ the distribution of the directed $i$NN graph in dimension $i$, we have that 
\begin{align*}
\P_2(A^2_n)=\E_3[\nu(A^2_n|\cdot)]\le \P_3(A^3_n), 
\end{align*}
which gives the result.
% Further let $E^i$ denote the edge set associated to $\Z^i$ for $i\in\{2,3\}$. 
% $\omega$ we will denote a configuration of open directed edges in $E^3$ and by
% $E^3_{\bar x}$ 
% $E^2_x$ denote the set of two-dimensional two-neighbor directed edge configurations $\omega_x$ of $x$. Correspondingly, 
% Let $\pi$ denote the projection of the edges from $\omega_{(x,n)}$ to the edges of $x$ in $\Z^2$, in words we forget all edges facing up or down and only keep all remaining edges which we consider as part the edge set of $\Z^2$. 
% We now define, in an algorithmic fashion, a probability kernel from $(E^3_{\bar x})$
% Let us start by  defining a probability kernel $\nu_x$ from $E^3_{\bar x}$ to probabilities on $E^2_x$. 
% \begin{enumerate}
% \item If $|\pi(\omega_{(x,0)})|=3$, pick two of the three directed edges uniformly at random and define the result as $\omega_x=F(\pi(\omega_{(x,0)}))$.
% \item If $|\pi(\omega_{(x,0)})|=2$, define $\omega_x=\pi(\omega_{(x,0)})$.
% \item If $|\pi(\omega_{(x,0)})|=1$, keep the edge in $\Z^2$ and follow the existing arrow upwards, consider $\omega_{(x,1)}$ and do the following:
% \begin{enumerate}
% \item If $|\pi(\omega_{(x,1)})|\ge 2$, pick one of the directed edges in $\pi(\omega_{(x,1)})\cap \pi(\omega_{(x,0)})$ uniformly at random and define the result as $\omega_x=F(\pi(\omega_{(x,0)}),\pi(\omega_{(x,1)}))$.
% \item If $|\pi(\omega_{(x,1)})|=1$ and $\pi(\omega_{(x,1)})\neq \pi(\omega_{(x,0)})$, define $\omega_x=\pi(\omega_{(x,0)})\cup\pi(\omega_{(x,1)})$.
% \item If $|\pi(\omega_{(x,1)})|=1$ and $\pi(\omega_{(x,1)})= \pi(\omega_{(x,0)})$, follow the existing arrow upwards and repeat from (a) with $\omega_{(x,1)}$ replaced by $\omega_{(x,2)}$.
% \end{enumerate}
% \end{enumerate}
% Now, almost surely, for every $x\in\Z^d$, the construction ends. Further, due to the complete symmetry in the construction, the resulting field $(\omega_x)_{x\in\Z^2}$ has the same distribution as the directed 2NN graph, i.e., $$\E_3[\nu(\d\omega|\cdot)]=\P_3(\d\omega).$$
% The essential benefit of the coupling comes from the following fact. Let  $\bar\omega$ be such that the origin is not connected to the boundary of a box with sidelength $n$, i.e., $\bar\omega\notin A^3_n$, then also $\nu(A^2_n|\bar\omega)=0$, since we are only using edges that are already present in $\bar\omega$. THIS I GUESS IS NOT QUITE CORRECT YET; WE HAVE TO DESIGN THE ALGORITHM STARTING FROM THE ORIGIN. 
% But then, 
% \begin{align*}
% \P_3(A^2_n)=\E_3[\nu(A^2_n|\cdot)]\le \P_3(A^3_n), 
% \end{align*}
% which tends to zero as $n$ tends to infinity, by assumption.
\end{proof}
\basc{I guess it could be quite annoying to write it out, but at least intuitively  we should be able to argue that this coupling works for any $d=k\geq 3$ in general, right?}



It would be desirable to proof the following conjecture.
\begin{conj}
$\theta^{\rm D}(2,2)>0$. 
\end{conj}
IDEA: USE PEIERLS ARGUMENT


\subsection{Results for the bidirectional $k$-neighbor graph}\label{sec:B}
Let us denote by $c(d)$ the connective constant of $\Z^d$, that is $$c(d)=\lim_{n\uparrow\infty}c_n(d)^{1/n},$$
where $c_n(d)$ is the number of self-avoiding paths in the $d$-dimensional cubic lattice, starting at the origin, of length $n$. The limit exists by subadditivity. Also, $d<c(d)<2d-1$.  
\begin{lemma}\label{lemma-Bnoperc}
For all $k$ such that $k(k-1)<2d(2d-1)/c(d)$ we have that $\theta^{\rm B}(k,d)=0$.
\end{lemma}
Let us note that this results can be interpreted in two ways. First, for given $d$, absence of percolation is guaranteed for sufficiently small $k$, and for example $\theta^{\rm B}(2,d)=0$ for any $d\geq 2$ (while it is clear that $\theta^{\rm B}(2,1)=1$). However, also for fixed $k$, since $2d(2d-1)/c(d)>2d$, for sufficiently large $d$, there is no percolation. This is due to the fact that in high dimensions, it is unlikely for two neighboring vertices to pick the connecting edge. 
\begin{proof}
Note first that the system is negatively correlated in the sense that for all directed edges $\ell_1=(x_1,x_2)\neq\ell_2=(y_1,y_2)$ we have that
$$
\mathbb P(\ell_1\text{ open and }\ell_2\text{ open})\le \mathbb P(\ell_1\text{ open})^2.
$$
Indeed, note that if $x_1\neq y_1$, then the inequality is an equality by independence. However, if $x_1=y_1$ then 
$$
\mathbb P(\ell_1\text{ open}|\ell_2\text{ open})=\frac{k-1}{2d-1}\le \frac{k}{2d}=\mathbb P(\ell_1\text{ open}).
$$
Using the first moment method, we have that 
\begin{align*}
\P(o\leftrightsquigarrow_{\rm B}\partial B_n)\le \sum_{s\in \mathcal R_n}\P(s\text{ is open})\le C (c(d)p)^n 
\end{align*}
where $p=k(k-1)/(2d(2d-1))$ is the probability that the origin chooses precisely two prescribed edges. This leads to the criterion for absence of percolation $c(d)p<1$.
\end{proof}

We need some result on nontrivial percolation behavior
\begin{prop}\label{prop-1indep}
\begin{enumerate}[(i)]
\item For any $d \geq 2$, if $k>2\sqrt{0.8457}d \approx 1.8392d$, then $\theta^{\rm B}(k,d)>0$.
\item For any $\alpha > 2 \sqrt{0.5847} \approx 1.5293$, we have $\theta^{\rm B}(\alpha d, d)>0$ for all $d$ sufficiently large.
\end{enumerate}
\end{prop}
\begin{proof}
Note that the $k$-BnG model (just as the $k$-UnG model) in any dimension $d$ is an 1-independent Bernoulli bond percolation model where each edge is open with probability $p=k^2/(4d^2)$. Let us write 
\[ p_{\max}(\Z^d)= \sup \{ p \in (0,1) \colon \text{some 1-independent Bernoulli bond perc.\ model does not percolate} \}. \]
According to [Balister et al.], we have $p_{\max}(\Z^d) \leq p_{\max}(\Z^{d-1})$ and $p_{\max}(\Z^2) \leq 0.8457$. These assertions imply (i). Part (ii) follows from the assertion of the same paper that $\lim_{d\to\infty} p_{\max}(\Z^d) \leq 0.8457$. 
\end{proof}
\begin{remark}
The smallest nontrivial (i.e., such that $k \neq 2d$) case covered by Proposition~\ref{prop-1indep} is $k=13$ and $d=7$: Since $13/7 >1.8392$, it follows that $\theta^{\rm B}(13,7)>0$. Similarly, $\theta^{\rm B}(24,13)>0$ is the first case where percolation follows for $k=2d-2$, $\theta^{\rm B}(35,19)$ is the first one for $k=2d-3$, and so on.

Although the $k$-UnG model is also an 1-independent Bernoulli bond percolation model, the same approach unfortunately does not yield any new results in its case. Indeed, here each edge is open with probability $p=k(4d-k)/(4d^2)$. Thus, for $k=d$ it always holds that $p=3/4<0.8457$, while for $k>d$ we already know that $\theta^{\rm D}(k,d)=1$. 
\end{remark}

%\subsection{Results for the exclusively unidirectional $k$-neighbor graph}
%Since the $k$-XnG is a subgraph of the $k$-UnG, the following assertion is immediate.
%\begin{prop}
%We have that $\theta^{\rm X}(1,d)=0$ for all $d \geq 1$. 
%\end{prop}



%Using a similar argument, we can also prove the followin
%\begin{lemma}
%We have $\theta(k,d)\leq \theta (k,d')$ whenever $k\geq 1$ and $d\leq d'$.
%\end{lemma}
%\begin{proof}
%The lemma follows from the fact that the restriction of the $k$-nnG in $d'$ dimensions to a $d$-dimensional coordinate plane is stochastically dominated by a $k$-nnG in $d$ dimensions.
%\end{proof}

%\newpage

%%OLD COX-DURRETT BOUND FOR UNDIRECTED/BIDIRECTIONAL, IT HAS ERRORS
%THIS DOESN'T SEEM TO WORK OUT DUE TO PROBLEMS WITH THE MARTINGALES.
%
%With methods analogous to the case of the $k$-UnG, we can additionally prove the following.
%\begin{prop}\label{prop-U}
 %   We have $\theta^{\rm U}(3,4)>0$. Moreover, $\theta^{\rm U}(2,d)>0$ for all $d \geq 5$. \color{red} Need to check issues with lack of independence. \color{black}
%\end{prop}
%\begin{proof}
%We argue analogously to the proof of Theorem~\ref{thm-CoxDurrett}, the only difference being that now the probability of an edge being present in the graph is higher. \color{red} (Add negativity of correlations from the old stuff. \color{black} Namely, we have for an undirected edge $\ell$ with endpoints $x_1,x_2 \in \Z^d$ being nearest neighbors:
%\begin{equation}\label{Uprob}
%\begin{aligned}
%\P&(\ell\text{ is an edge of the $k$-UnG}) = \P((x_1,x_2)\text{ is an edge of the $k$-UnG}) \\ & \quad +\P((x_2,x_1)\text{ is an edge of the $k$-UnG})- \P(\text{both } (x_1,x_2) \text{ and } (x_2,x_1) \text{ are edges of the $k$-UnG} ) \\ &= \frac{k}{2d} + \frac{k}{2d} - \frac{k^2}{4d^2} = \frac{k}{d} \Big( 1 - \frac{k}{4d} \Big),
%\end{aligned}
%\end{equation}
%since the events $\{ (x_1,x_2)\text{ is an edge of the $k$-UnG}\}$ and $\{ (x_2,x_1) \text{ is an edge of the $k$-UnG} \}$ are independent and their probabilities equal $k/(2d)$. Now, using the martingale $W_n=N_n/((\frac{k}{d} ( 1 - \frac{k}{4d} ))^n)$, if $\varrho(d)< \frac{k}{d} ( 1 - \frac{k}{4d} )$, it follows analogously to the proof of Theorem~\ref{thm-CoxDurrett} that $\theta^{\rm U}(k,d)>0$. 
%
%Let us first consider the case $(k,d)=(3,4)$. In the proof of Theorem~\ref{thm-CoxDurrett} we computed that $\varrho(4)=\P(\tau_4<\infty) < 0.5$, and thus
%\[ \varrho(4) < \frac{k}{d} \Big( 1 - \frac{k}{4d} \Big)= \frac{3}{4} \Big( 1- \frac{3}{4 \cdot 4} \Big) = \frac{39}{64} = 0.609375. \] 
%Thus, $\theta^{\rm U}(3,4)>0$.
%
%Let us now focus on the case $(k,d)=(2,5)$. We also computed in the aforementioned proof that $\varrho(5)<0.28$, and thus
%\[ \varrho(5) < \frac{k}{d} \Big( 1 - \frac{k}{4d} \Big) = \frac{2}{5} \Big( 1- \frac{2}{20} \Big) = 0.36. \]
%This implies that $\theta^{\rm U}(2,5)>0$. For $d=6,7,8,9,10$ it in fact holds that the right-hand side of~\eqref{CDUB} is smaller than $\frac{2}{d} (1-2/(4d))$, and for $d \geq 10$ it is again elementary to show that this right-hand side of~\eqref{CDUB} minus $\frac{2}{d} (1-2/(4d))$ stays negative (given that it is negative for $d=9$). \color{red} (Indeed, I think that we can again argue that the rhs of~\eqref{CDUB} for dimension $d+1$ is at most $d/(d+1)$ times the same for dimension $d$, while $\frac{2}{d+1}(1-2/(4d+1))>\frac{d}{d+1} \frac{2}{d}(1-2/(4d))$. \color{black}) Therefore, $\theta^{\rm U}(2,k)>0$ holds for all $k\geq 5$, as asserted.
%\end{proof}
%
%\color{green} If this is fine, the conjecture is proven apart from $(k,d)=(2,2),(2,3),(3,3),(2,4)$. For $(2,4)$ it is not impossible that a Cox--Durrett style bound can work, but we would need $\P(\tau_4=k)$ at least for all $k\leq 8$. So the problem is similar as in the directed case for $k=3$ and $d=4$. \color{black}
%
%
%\begin{prop}\label{prop-B}
%For $d \geq 5$ and $k \geq \sqrt{6 d}$ and for $d \geq 4$ and $k\geq \sqrt{8d}$ we have $\theta^{\rm B}(k,d)>0$.
%\end{prop}
%This implies e.g.\ that $\theta^{\rm B}(5,5)>0$ and $\theta^{\rm B}(6,4)>0$. \color{green} So far, $k=5$ is the smallest $k$ for which we get percolation for some $d$. \color{black}
%\begin{proof}[Proof of Proposition~\ref{prop-B}]
%We again adapt the proof of Theorem~\ref{thm-CoxDurrett}, similarly to how we did it in the proof of Proposition~\ref{prop-U}. \color{red} Also here we have negative correlations -- need to copy it from old stuff. \color{black} For an undirected edge $\ell$ with endpoints $x_1,x_2 \in \Z^d$ being nearest neighbors we have (cf.~\eqref{Uprob})
%\[ \P(\ell \text{ is an edge of the $k$-BnG}) = \P(\text{both } (x_1,x_2) \text{ and } (x_2,x_1) \text{ are edges of the $k$-UnG} )=\frac{k^2}{4d^2}. \]
%Now, we have seen in the aforementioned proofs that $\varrho(d)<\frac{3}{2d}$ for $d \geq 5$ and $\varrho(d)<\frac{2}{d}$ for $d \geq 4$. Consequently, using the martingale $W_n=(4d^2/k^2)^n N_n$, we obtain percolation for $d\geq 5$ whenever 
%\[ \frac{3}{2d} < \frac{k^2}{4d^2}, \]
%which is equivalent to $k \geq \sqrt{6 d}$. For $d \geq 4$ we obtain percolation whenever
%\[ \frac{2}{d} < \frac{k^2}{4d^2}, \]
%in other words, if $k \geq \sqrt{8d}$, as claimed.
%\end{proof}


\subsection{Results for the undirected $k$-neighbor graph}\label{sec:U}
We need some results here.
\begin{conj}
We have that $\theta^{\rm U}(k,d)>0$ for all $k,d\ge 2$.  
\end{conj}

\medskip

It is clear that $\theta^{\rm U}(k,d) \geq \theta^{\rm D}(k,d)$ for any $k,d$ because any directed edge in the $k$-DnG corresponds to an undirected edge in the $k$-UnG (and at most two directed edges can correspond to the same undirected edge). Additionally, $k=1$ is not sufficient for percolation even in the undirected sense.

\begin{prop}\label{prop-U1}
    We have that $\theta^{\rm U}(1,d)=0$ for all $d \geq 1$.
\end{prop}
\begin{proof}
Assume for a contradiction that the assertion of the proposition is wrong, and let $A$ then denote the event of positive probability that there exists an infinite path starting from $o$ consisting of edges of the $1$-UnG. %First, we claim that if such a path exists, it can contain at most one bidirectional edge. Indeed, two bidirectional edges cannot be adjacent to each other because this would mean that the vertex separating them has outdegree at least 2 in the $1$-DnG. Also, vertices belonging to the bidirectional edge have no outgoing edge but the one contributing to the bidirectional edge. This way, if $Y_0=o, Y_1,Y_2,\ldots$ is the enumeration of the vertices among an infinite path and for some $i \geq 0$, $(Y_i,Y_{i+1})$ is a bidirectional edge, then either $Y_{i+2}=Y_i$ and $Y_i$ has an additional incoming edge in the $1$-DnG coming from $Y_{i+3} \neq Y_{i+1}$, or $Y_{i+2} \neq Y_i$ and $Y_{i+1}$ has an additional incoming edge coming from $Y_{i+2}$. Now, 
%such a path, then if there is a bidirectional edge, say $(Y_{i},Y_{i+1})$, then the only way that this edge is in the infinite path is that $Y_i$ or $Y_{i+1}$ has an additional incoming edge from some other vertex
%there also exists one with the same property that does not contain any bidirectional edge apart from the very first one (which is then adjacent to $o$). Indeed, in any infinite path of the $1$-UnG, there cannot be two bidirectional edges adjacent to each other because this would mean that the vertex between them has outdegree at least 2. Thus, taking any infinite path starting from $o$, this can only contain isolated bidirectional edges (i.e., ones not adjacent to any other one). Hence, if $Y_0=o,Y_1,Y_2,\ldots$ is an infinite path of the $1$-UnG and for some $k \geq 0$, $(Y_k,Y_{k+1})$ is a bidirectional edge, it must be the case that either $k=0$, or $(Y_{k-1},Y_k)$ and $(Y_{k+2},Y_{k+1})$ is an edge of the $1$-DnG. In either of the cases, if we leave out the bidirectional edge $(Y_k,Y_{k+1})$ from the infinite path, the path still remains infinite, and this can be done simultaneously with all bidirectional edges of the path. This implies the claim.
let $X_0=o,X_1,X_2,\ldots$ be any infinite path of the $1$-UnG starting from $o$.
%that contains no bidirectional edges apart from $(X_0,X_1)$ and is defined (on the event $A$) as follows. First, consider the infinite cluster of $o$ in the $1$-UnG and remove all bidirectional edges of it that are not adjacent to $o$, along with the vertices that are only adjacent to a bidirectional edge of the cluster. Now, we define $X_1,X_2,\ldots$ recursively as follows. If $o$ is not adjacent to a bidirectional edge of the infinite cluster, then given that $X_0,\ldots,X_{i-1}$ is already defined, let $X_i$ be the lexicographically smallest one among all vertices $v$ of the remainder of the infinite cluster of $o$ such that $(X_{i-1},v)$ is an edge of the $1$-UnG and there is a path in the $1$-UnG from $v$ to infinity that does not use any of the edges $(X_0,X_1),\ldots,(X_{i-1},v)$. (Note that there may be several $v$ satisfying this property in case there are multiple edge-disjoint incoming paths of the $1$-DnG at $X_{i-1}$.) If $o$ is adjacent to a bidirectional edge, then let $X_1$ be its neighbor via this edge, and define $X_2,X_3,\ldots$ as before.

Now, if $X_0=o,X_1,X_2$ is the sequence of vertices in an infinite path (on the event $A$), we put
\[ K=\inf \{ k \geq 0 \colon (X_k,X_{k+1}) \text{ is not an edge of the $1$-DnG} \}. \]
We know from (the proof of) Lemma~\ref{lem:surenoperc} that $K<\infty$ a.s. Next, let us define
\[ L=\inf \{ k > K \colon 
(X_{k},X_{k+1})\text{ is an edge of the 1-DnG}\}. \]
Now we claim that $L=\infty$ on the event $\{ K<\infty\}$ (where we put $K=\infty$ on the event $A^c$). Indeed, for all $k=K,K+1,\ldots,L-1$, $(X_{k+1},X_k)$ is an edge of the $1$-DnG because it follows from the definition of $K$ that $(X_{k},X_{k+1})$ is not an edge of the 1-DnG, but $(X_k,X_{k+1})$ is an edge of the 1-UnG, and this is only possible if $(X_{k+1},X_k)$ is an edge of the $1$-DnG. But now, if $L<\infty$, then $(X_{L},X_{L-1})$ is an edge of the $1$-DnG, but then $(X_{L},X_{L+1})$ cannot be an edge of the $1$-DnG since there is only one edge going out of $X_L$. This implies the claim.

%\color{blue} Note that this implies that the cluster of $o$ in the $1$-UnG consists of coalescing (finite or infinite) directed paths towards $X_K$ and $X_{K+1}$ if $K\geq 1$; if $K=0$, then all are directed towards $o$. In particular, there can be at most one bidirectional edge in the infinite cluster, namely the one between $X_{K-1}$ and $X_K$ (because if two vertices have outgoing edges of the $1$-DnG towards each other, they cannot have any other outgoing edges, so that $K$ is at most the larger one among the indices of the two endpoints of the bidirectional edge). We perhaps don't need this, but I wanted to note it. \color{black}

Let us denote the $\ell^1$-sphere of radius $n\in\N_0$ by $S^{(1)}_n$ (with $S^{(1)}_0=\{o\}$). From Lemma~\ref{lem:surenoperc} we even know that $\P(\kappa<\infty\,|\,A)$? where
\[ \kappa = \inf \big\{ k \geq 0 \colon o \not\rightsquigarrow S^{(1)}_k \text{ in the $1$-DnG} \big\}. \]
By continuity of measures, we can find $k_0\in\N$ such that $\P(\kappa \leq k_0 |A) \geq 1/2$. Now, by our previous observation, since from $o$ to $S^{(1)}_\kappa$ there always exists a directed path in the $1$-DnG, 
on the event $\{ \kappa \leq k_0 \}$, for all $n\in\N_0$ there must exist a self-avoiding directed path of length $n$ in the $1$-DnG ending at some vertex of $S^{(1)}_{k_0}$. Now, writing $c_{n,k}(d)$ for the number of self-avoiding paths in the $d$-dimensional lattice of length $n$ starting from (or ending at) $S^{(1)}_{k_0}$, it is clear that
\[ \lim_{n\to\infty} c_{n,k}(d)^{1/n} = \lim_{n\to\infty} c_{n,0}(d)^{1/n}=c(d) \leq 2d-1, \]
where $c(d)$ is the connective constant of $\Z^d$ (cf.\ Section~\ref{sec:B}). Now, any self-avoiding path of length $n$ ending at $S^{(1)}_{k_0}$ is open with probability $(2d)^{-n}$. Hence,
\[ \P(\exists\text{ a self-avoiding path of length $n$ starting from $S^{(1)}_{k_0}$ included in the $1$-DnG}) \leq \Big( \frac{2d-1}{2d} \Big)^n, \]
which decays exponentially fast in $n$. This contradicts the assertion that the left-hand side is bounded from below by $\P(A)/2>0$ independently of $n$. Therefore, $\P(A)=0$, as wanted.
\end{proof}

Next, we verify percolation in the $2$-UnG model in two dimensions.

\begin{theorem}\label{theorem-U22}
  We have $\theta^{\rm U}(2,2)>0$. 
\end{theorem}

\begin{proof}
    We prove the positive percolation probability via a dual approach. That is, let $\D^2:=\{(x+\frac12,\frac12): x\in \Z^2\}$ be the two-dimensional dual lattice of $\Z^2$. An edge $e$ in the dual lattice has exactly one edge $e'$ from the lattice $\Z^2$ that crosses $e$, and we say the edge $e$ is closed if and only if the edge $e'$ is open. For the $2-UnG$ model in two dimensions, it thus follows that for any such edge $e$, 
    \be 
    \P(e \text{ is closed})=1-\P(e'\text{ is closed})=1-\frac34=\frac14.
    \ee 
    Moreover, the negative correlations between edge that are present in the $2$-UnG model also appear in its dual. Namely, for two edges in the dual that are vertex disjoint or not orthogonal, the status of the edges is independent. When two edges $e_1,e_2$ in the dual share a vertex and are orthogonal, that is, when there is some vertex $v\in \Z^2$ such that the edges emanating from $v$ can cross both $e_1$ and $e_2$, then
    \be 
    \P(e_1\text{ and }e_2\text{ closed in the dual})=\frac{1}{24}\leq \Big(\frac14\Big)^2=\P(e_1\text{ closed in the dual})\P(e_2\text{ closed in the dual}).
    \ee 
    This can be extended, so that for any $k\geq 2$ and any path $(v_1,\ldots, v_k)$ of nearest-neighbour vertices in $\D^2$, 
    \be 
    \P(\text{Path }(v_1,\ldots, v_k)\text{ closed})\leq \prod_{i=1}^{k-1}\P{(v_i,v_{i+1})\text{ closed}}=\Big(\frac14\Big)^{k-1}.
    \ee 
    As such, we can use the negative correlation to bound the probability that there exists a closed circuit in the dual that surrounds the origin. For any $n\geq 4$, the number of such circuits of length $n$ is at most $n\sigma(n-1)$, where $\sigma(n-1)$ denotes the number of self-avoiding paths in $\Z^2$. As $\sigma(n-1)=(c(d)+o(1))^{n-1}$, we thus arrive at 
    \be \ba\label{eq:finnorpercprob}
    1-\theta^U(2,2)&=\P(\exists \text{ circuit }\gamma \text{ in }\D^2\text{ that surrounds the origin})\\
    &\leq \sum_{n=4}^\infty \sum_{\substack{\text{circuits }\gamma\\ |\gamma|=n}}\P(\gamma\text{ closed in }\D^2)\\
    &\leq \sum_{n=4}^\infty n(c(2)+o(1))^{n-1} \Big(\frac14\Big)^n, 
    \ea \ee 
    which is finite since $c(2)\leq 3$. We now apply an argument from {\color{red} Cite Durret 88 or Jan Swart's notes} to show that this probability is in fact smaller than one, which yields $\theta^U(2,2)>0$. Let $D_m=\{0,\ldots, m\}^2$ and say that a vertex $i\in \Z^2$ is `wet' when there exists a $j\in D_m$ such that $i\to j$, that is, when there exists a path of closed edges from $i$ to $j$. Suppose the component of the origin is finite almost surely. Then, the number of wet vertices is finite almost surely as well, and there exists a circuit in the dual that surrounds all the wet sites. By the same argument as in~\eqref{eq:finnorpercprob}, 
    \be 
    \P(D_m \not\to \infty )\leq \sum_{n=4m}^\infty \sum_{\substack{\text{circuits }\gamma \\ |\gamma|=n}}\P(\gamma\text{ closed in }\D^2)\leq \sum_{n=4m}^\infty n(c(d)+o(1))\Big(\frac14\Big)^n,
    \ee 
    where we can start the outer sum from $4m$ since a circuit that surrounds all wet vertices must surround $D_m$ and hence have length $4m$ at least.\basc{Do we need to allow only vertices in $\Z^2\backslash D_m$ to be wet? Otherwise, if vertices in $D_m$ can be wet (and in particular, it can then be possible that \emph{only} vertices in $D_m$ are wet). If that would be the case, I don't see why such a circuit must surround $D_m$.} 
    As a result, since the sum is finite for each $m\in\N$, choosing $m$ large enough yields that $\P(D_m\to \infty)>0$ and hence $\P(i\to \infty)>0$ for some $i\in D_m$. By translation invariance, it then follows that $\theta^U(2,2)>0$, as desired.
\end{proof}



\section{Discussion}\label{sec:discussion}
\subsection{Phase transitions from the perspective of complementary relations}\label{sec:phases}
Let the graphs $G,G_1,G_2$ all be directed (resp.\ all undirected) with the same vertex set. We say that $G_1$ is the complement of $G_2$ w.r.t.\ $G$ if the edge sets of $G_1$ and $G_2$ are disjoint and their union equals the edge set of $G$. For undirected graphs with $n \in \N$ edges, choosing $G=K_n$ yields the usual definition of graph-theoretic complement. Then, the following assertions hold for $1 \leq k < 2d$.
\begin{itemize}
    \item The complement of the $k$-DnG with respect to $\Z^d$ with all directed nearest neighbor edges is a $2d-k$-DnG. (In particular, the $d$-DnG equals its such complement in distribution.)
    \item The complement of the $k$-UnG with respect to $\Z^d$ with all undirected nearest neighbor edges is a $2d-k$-BnG. (Hence, the same is also true with the words ,,UnG'' and ,,BnG'' swapped.) 
    \item $\Z^d$ with all directed nearest neighbor edges is a $2d$-DnG and with all undirected nearest neighbor edges is a $2d$-UnG and a $2d$-BnG at the same time. Hence, the above two assertions also hold for $k \in \{ 0, 2d\}$, where the 0-DnG is defined as $\Z^d$ as a directed graph with no edges and the 0-UnG and 0-BnG as $\Z^d$ as an undirected graph without edges.
    %\item For $1 \leq k \leq 2d$, the graph-theoretic complement of the $k$-XnG is the disjoint union of two correlated graphs, their marginal distributions being the one of a $k$-BnG and a $2d-k$-BnG. 
\end{itemize}

In what follows, we will simply speak about complements of the graphs w.r.t.\ $\Z^d$, which should always be understood in the corresponding (directed resp.\ undirected) sense. Let now $d \in \N$ and $1 \leq k < 2d$. Borrowing a nomenclature known from continuum percolation \cite[Section 7.3]{MeesterRoy1996}, we say that 
\begin{itemize}
    \item the $k$-UnG in $d$ dimensions is an \emph{infinite cascade} if $\theta^{\rm U}(k,d)=0$ and $\theta^{\rm B}(2d-k,d)=0$,
    \item the $k$-UnG in $d$ dimensions is \emph{supercritical with subcritical complement} if $\theta^{\rm U}(k,d)>0$ and $\theta^{\rm B}(2d-k,d)=0$,
    \item the $k$-UnG in $d$ dimensions is \emph{subcritical with supercritical complement} if $\theta^{\rm U}(k,d)>0$ and $\theta^{\rm B}(2d-k,d)=0$,
    \item the $k$-UnG in $d$ dimensions exhibits a \emph{coexistence of phases} $\theta^{\rm U}(k,d)>0$ and $\theta^{\rm B}(2d-k,d)>0$,
\end{itemize}
The same notions can be defined analogously for the $k$-BnG, swapping the roles of $\rm U$ and $\rm B$ everywhere in the above definitions, and for the $k$-DnG, replacing each $\rm U$ and $\rm B$ by a $\rm D$. In words, an infinite cascade means that neither the graph nor its complement (w.r.t.\ $\Z^d$) has an infinite connected component, supercriticality with subcritical complement means that the graph has one but its complement has none, while a coexistence of phases means that both graphs have an infinite cluster. \color{red} In many percolation models, it is the case that a coexistence of phases requires a high dimension (typically higher than 2) -- references? \color{black}

Now, for $k=d$, since the $k$-DnG is equal to its own complement w.r.t.\ $\Z^d$ in distribution, either it is an infinite cascade or it exhibits a coexistence of phases. Further, for $k=d$, the complement of the $k$-BnG w.r.t.\ $\Z^d$ equals the $k$-UnG in distribution, and hence the $k$-BnG cannot be supercritical with subcritical complement, while the $k$-UnG cannot be subcritical with supercritical complement. Finally, if the $k$-DnG is supercritical (resp.\ subcritical) with subcritical (resp.\ supercritical) complement, then $k>d$ (resp.\ $k<d$). In the directed case, the fact that $\theta^{\rm D}(k,d)=1$ for $k>d$ implies that an infinite cascade is only possible for $k=d$.

According to our main results, the following assertions hold.
\begin{itemize}
    \item It is trivial that in one dimension, the $1$-DnG is an infinite cascade, and so are the $1$-UnG and the $1$-BnG. Even if there are nontrivial examples for an infinite cascade, there are only finitely many of them. The only possible examples in the directed case are $k=d=2$ and $k=d=3$. However, Theorem~\ref{theorem-3322}  implies that if the $3$-DnG for $d=3$ is an infinite cascade, then so is the $2$-DnG for $d=2$, and as already anticipated, numerical simulations actually suggest that in this case we have a coexistence of phases, see \color{red} ??? Bas' figure? \color{black} . In the undirected cases, the possible cases are $2 \leq d \leq 6$ for $k=1$, $3 \leq d \leq 12$ for $k=2$ and $3 \leq d \leq 4$ for $k=3$.
    \item Supercriticality with subcritical complement occurs in all dimensions, both in the directed and undirected cases. Since the $k$-DnG (or the $k$-UnG) percolates for $k\geq 3$ if $d \geq 5$, in these dimensions it can only be subcritical with supercritical component for $k=1,2$, and we know that for $k=1$ this is actually the case, even for $d \geq 2$. The ``sporadic'' open case $k=3, d=4$ that cannot be an infinite cascade could be another example where the $k$-DnG is subcritical with supercritical complement; otherwise we have a coexistence of phases. Also for $k=3$ and $d=4$, but for the UnG, none of the four possibilities can be excluded. For $d=2$, we have proven that the $2$-UnG is supercritical with subcritical complement. For $k \geq 1$ fixed, we have seen that the same holds for the $2d-k$-UnG if $d$ is sufficiently large (and for $k=1$ actually for all $d \geq 2$ and for $k=2$ for all $d \geq 3$). 
    \item A coexistence of phases occurs in many high-dimensional cases both in the directed and the undirected--bidirectional model. In the directed case, for $d \geq 5$, since $2d-3 > d$ for such a dimension, it follows that there is a coexistence of phases for $k \in \{3, \ldots, 2d-3 \}$.  We also observe a coexistence of phases for $k=d=4$, which is the smallest-dimensional known example of a coexistence of phases so far. Also for fixed $k \geq 3$, there is a coexistence of phases for the $k$-UnG for $d$ sufficiently large, and the same holds in general for $k=k(d)$ such that $k(d) < 2d-2 \sqrt{0.5847}d$. The smallest-dimensional example with coexistence of phases in the undirected--bidirectional model is the one of the $3$-UnG (and the $35$-BnG) in 19 dimensions.  It would be interesting to know if a coexistence of phases can be observed in any of our models for $d=2,3$; a natural candidate mentioned already is the $2$-DnG in 2 dimensions.
\end{itemize}
In this discussion we have not distinguished between the cases when $o \rightsquigarrow \infty$ happens with probability in $(0,1)$ or with probability 1. The latter is the case in the $k$-DnG and $k$-UnG if (and only if!) $k>d$ (cf.\ Lemma~\ref{lem:sureperc}). This does not exclude a coexistence of phases in the directed case (take e.g.\ $k=3$ and $d \geq 5$). It is an interesting open question whether the same holds for the undirected--bidirectional case. We are not aware of cases where $k<d$ and $\theta^{\rm B}(k,d)>0$, which would be necessary for such an assertion.

\appendix

\section{Proof of Lemma~\ref{lemma-largedmon}}\label{sec:largedmon}
In this section, we carry out the proof of Lemma~\ref{lemma-largedmon}.
\begin{proof}[Proof of Lemma~\ref{lemma-largedmon}]
Recall that~\eqref{largedmon}, i.e.,
\[ R(k,d):=\frac{1}{d} + \frac{1}{d^3} -\frac{1}{d^4} + \sum_{k=3}^d d^{-k} k! +  (2\pi d)^{1/2} \Big( \frac{\e^{-1/13}}{\sqrt{2\pi}} \Big)^d \zeta\Big( \frac{d-1}{2} \Big) - \frac{k}{2d} < 0.  \]
holds for any $k\geq 3$ if $d=7$. 
Our goal is to prove the statement of Lemma~\ref{lemma-largedmon}, i.e.\ that the same holds for any $d \geq 7$ and $k\geq 3$. We check numerically that it also holds for $8 \leq d \leq 11$ and $k=3$ (and thus also for $k\geq 4$ for the same values of $d$): indeed, we have $R(8,3)=-0.0292277$, $R(9,3)=-0.0350912$, $R(10,3)=-0.0367514$ and $R(11,3)=-0.0364418$. 

Now, it is easy to show that for $d \geq 11$, we have
\begin{equation}\label{Rfactor} R(k,d+1) \leq \frac{d}{d+1} R(k,d) + (d+1)^{-(d+1)} (d+1)! \end{equation}
Indeed, it holds that
\[ 
\begin{aligned}
   & R(k,d)-(d+1)^{-(d+1)}(d+1)! \\ &=\frac{1}{d+1}- \frac{k}{2(d+1)} + \frac{1}{(d+1)^3} -\frac{1}{(d+1)^4} + \sum_{k=3}^{d} (d+1)^{-k} k! +  (2\pi (d+1)^{1/2} \Big( \frac{\e^{-1/13}}{\sqrt{2\pi}} \Big)^{d+1} \zeta\Big( \frac{d}{2} \Big)
\\    &=\frac{d}{d+1} \Big( \frac{1}{d}- \frac{k}{2d} \Big) + \frac{1}{(d+1)^3} -\frac{1}{(d+1)^4} + \sum_{k=3}^{d} (d+1)^{-k} k! +  (2\pi (d+1)^{1/2} \Big( \frac{\e^{-1/13}}{\sqrt{2\pi}} \Big)^{d+1} \zeta\Big( \frac{d}{2} \Big)
\\ & \leq \frac{d}{d+1} \Big( \frac{1}{d}- \frac{k}{2d} \Big) + \frac{d}{d+1} \Big( \frac{1}{d^3} -\frac{1}{d^4} + \sum_{k=3}^d d^{-k} k! +  (2\pi d)^{1/2} \Big( \frac{\e^{-1/13}}{\sqrt{2\pi}} \Big)^d \zeta\Big( \frac{d-1}{2} \Big) \Big) \\ &=\frac{d}{d+1} R(k,d),
\end{aligned}  \]
where in the last step we used that $\zeta$ is monotone decreasing in $(1,\infty)$ and that for $d \geq 11$,
\begin{equation*} \sqrt{\frac{2\pi(d+1)}{2\pi d}} \frac{\e^{-1/13}}{\sqrt{2\pi}} \leq \sqrt{\frac{12}{11}} \frac{\e^{-1/13}}{\sqrt{2\pi}}  \approx 0.385831, \end{equation*}
while $d/(d+1)>11/12 = 0.91\dot{6}$. 

Next, let us show that
\begin{equation}\label{newtermbound} (d+1)^{-(d+1)} (d+1)! \leq \frac{d}{d+1} \big( \frac{1}{d^3}-\frac{1}{d^4} \big) - \Big( \frac{1}{(d+1)^3} - \frac{1}{(d+1)^4} \Big) \end{equation}
holds for $d \geq 11$. Now note that for such $d$, we have
\[ (d+1)^{-(d+1)} (d+1)! \leq \frac{12!}{(d+1)^{12}} \leq \frac{12!}{12^8} \frac{1}{(d+1)^4} \approx \frac{1.114}{(d+1)^4} \]
and
\[ \frac{d}{d+1} \big( \frac{1}{d^3}-\frac{1}{d^4} \big) - \big( \frac{1}{(d+1)^3} - \frac{1}{(d+1)^4} \big) = \frac{(d-1)(d+1)^3}{d^3(d+1)^4} = \frac{2d^3-2d-1}{d^3(d+1)^4} > \frac{1.99}{(d+1)^4} \]
Hence,~\eqref{newtermbound} holds for $d \geq 11$, and thus such $d$ and for any $k\geq 3$,
\[ R(k,d+1) \leq \frac{d}{d+1} R(k,d). \]
Since $R(k,11) < 0$, the lemma now follows by induction over $d$.
\end{proof}

\section*{Acknowledgements}
The authors thank D.~Mitsche, G.~Pete and B.~Ráth for interesting discussions and comments.

\printbibliography



\newpage
NOT USEFUL, AT LEAST IN THIS CONTEXT:

\medskip
However, when we consider the undirected model in which the directed edges are replaced by undirected edges, we still observe negative correlations, 
$$
\mathbb P(\ell_1\text{ open}|\ell_2\text{ open})=\frac{k-1}{2d-1}+(1-\frac{k-1}{2d-1})\frac{k}{2d}\le \frac{k}{2d}(2-\frac{k}{2d})=\mathbb P(\ell_1\text{ open}).
$$
Similarly, for the bidirectional case
$$
\mathbb P(\ell_1\text{ open}|\ell_2\text{ open})=\frac{k-1}{2d-1}\frac{k}{2d}\le\frac{k^2}{4d^2}=\mathbb P(\ell_1\text{ open}).
$$
In this case, the first moment methods leads to the criterion
$$
k<\frac{2d}{\sqrt{c(d)}}
$$
and thus for example for $k=2$ the process is subcritical for any $d\ge 2$. For example for $k=3$ the process is subcritical for $d\ge 4$, or for $k=4$ the process is subcritical for $d\ge 8$.

Let us note that there is a major qualitative difference between the bidirection graph and the both other graphs. For the bidirectional graph it becomes harder to percolate in high dimensions since neighboring vertices find it harder to pick each other. This seems to be the dominate effect that overcomes the fact that also in higher dimensions it is harder to backtrack or do loops. On the other hand, without the need to make matching picks, the other graphs should percolate easier in higher dimensions since there only the backtracking and looping becomes less likely. 

In fact, we can be better since, in the bidirectional case, if the edges $\ell_1,\ldots,\ell_m$ form a path for $m \geq 2$, we have
$$
\mathbb P(\ell_1\text{ open}|(\ell_i)_{i=2}^{m-1} \text{ open})\le \mathbb P(\ell_1\text{ open}|\ell_2\text{ open})=\frac{k-1}{2d-1}\frac{k}{2d}
$$
and hence the criterion is
$$
\frac{k-1}{2d-1}\frac{k}{2d}c(d)<1.
$$
Since $c(d)<2d-1$, we have $k(k-1)\le 2d$, we have that for $k=3$ and $d=3$, the system is also subcritical. Hence we have the following theorem.
\begin{theorem}
Consider the bidirectional model with $k=3$. Then for $d=2$ the model is supercritical and for $d\ge 3$ the model is subcritical. 
\end{theorem}

Let us note that there are some fundamental differences between the lattice model and the continuum model in the bidirectional case. This can be seen for example by considering the expected degree of the origin. In the discrete case we have
$$
\mathbb E[N_k]=k\frac{k}{2d}, 
$$
which can be seen as the expected value of a Bernoulli experiment in which for each of the $k$ outgoing edges of the origin, we flip a coin to see if this edge is also an ingoing edge. In particular, this expected value tends to zero as $d$ tends to infinity. On the other hand in the continuum by the FKG inequality,
\begin{align*}
\mathbb E[N_k]&=\int_{\mathbb R^d}d x\mathbb P(X(B_{|x|}(o))\le k-1, X(B_{|x|}(x))\le k-1)\\
&\ge \int_{\mathbb R^d}d x\mathbb P(X(B_{|x|}(o))\le k-1)^2\\
&= \int_{\mathbb R^d}d x {\rm Pois}_{v_d|x|^d}(\{0,\dots,k-1\})^2\\
&= \int_0^\infty d r {\rm Pois}_{r}(\{0,\dots,k-1\})^2\\
&= \sum_{i,j=0}^{k-1}\frac{(i+j)!}{i!j!}2^{-1-i-j}
\end{align*}
where we also used the representation of the gamma function. This expression is independent of the dimension and for $k=2$ we have $\mathbb E[N_k]\ge \color{red} 5/4 \color{black}$. Since $\mathbb E[N_k]$ is increasing in $k$, we can not use the dimension to reduce the expected degree. Actually, in the limit $d$ to infinity, the FKG inequality becomes an equality. Also, from simulations, it should be true that $\lim_{k\uparrow\infty}\mathbb E[N_k]/k=1$, so in high dimensions the bidirectional model behaves more like the directional model where  $\mathbb P(N_k=k)=1$. It is also immediate that for the undirected case we have $\mathbb E[N_k]\le 2k$. What is the role of dimension in the continuum model?

For the record, the computation can be made precise without FKG, for the bidirectional model,
\begin{align*}
\mathbb E[N_k]&=\sum_{i=0}^{k-1}\sum_{j=0}^{k-1-i}\sum_{l=0}^{k-1-i}\frac{(i+j+l)!}{i!j!l!}\frac{\hat v_d^i\bar v_d^{j+l}}{(\hat v_d+2\bar v_d)^{1+i+j+l}}
\end{align*}
where $\hat v_d=|B^d_1(o)\cap B^d_1(e_1)|/|B^d_1(o)|$ and $\bar v_d=|B^d_1(o)\setminus B^d_1(e_1)|/|B^d_1(o)|$.

Let $u\sim v$ denote that two vertices $u,v\in \Z^d$ are nearest-neighbours and let $v\to u$ denote that $v$ has an outgoing edge directed towards $u$ in $k$-nn$G$ on $\Z^d$. 

\begin{lemma}[Expected number of neighbours further away from the origin]
	Fix $k\leq d$ and let $v\in \Z^d$ be a vertex such that $|\{u\in \Z^d: u\sim v\text{ and } \|u\|_1=\|v\|_1+1\}|=d$. Then, 
	\be 
	\E{|\{u\in \Z^d: v\rightarrow u\text{ and } \|u\|_1=\|v\|_1+1\}|}=\frac k2.
	\ee 
\end{lemma}

\begin{proof}
	As $v$ has exactly $d$ many nearest-neighbours such that $\|u\|_1=\|v\|_1+1$, we can write the expected value as 
	\be
	\sum_{j=0}^k j\frac{\binom{d}{j}\binom{d}{k-j}}{\binom{2d}{k}}=\frac{d}{\binom{2d}{k}}\sum_{j=0}^k \binom{d-1}{j-1}\binom{d}{k-j}=\frac{d}{\binom{2d}{k}}\sum_{j=0}^{k-1}\binom{d-1}{j}\binom{d}{(k-1)-j}.
	\ee 
	We then use the Vandermonde identity to arrive at 
	\be 
	\E{|\{u\in \Z^d: v\rightarrow u\text{ and } \|u\|_1=\|v\|_1+1\}|}=\frac{d\binom{2d-1}{k-1}}{\binom{2d}{k}}=\frac k2,
	\ee 
	as desired.
\end{proof} 

\subsection{Some ideas for a proof} \color{red} Is this useful or can it go away? -- András \color{black}

We consider the growths process starting at the origin $X^0=\{o\}$, $X^1=\{X_1,X_2\}$, where $X_1, X_2$ are the two neighbors of the origin, the two (out of the $4$) neighbors such that there is a directed edge from the origin to that neighbor. In the same vain we define $X^1=\{X_{1,1},X_{1,2},X_{2,1},X_{2,2}\}$ and  $X^n=\{X_{i_1,\dots,i_n}\}$ and note that in this construction we track the path from the origin to that point via the indices. However these paths are not self-avoiding and the set $X^n$ may contain the same point multiple times. Therefor we also consider the sets of newly discovered vertices
$$
Y^n=X^n\setminus X^{n-1} 
$$
and note that of course also $Y^n$ may contain points with different indices that are actual the same. Note that by construction, if the cluster of the origin is finite, then there must exists a generation $n$ with $Y^n=\emptyset$. 

Our starting point is the following standard ansatz
\begin{align*}
1-\theta(2,2)&= \mathbb P(\text{there exists $n$ such that $Y^n=\emptyset$})\\
&\le  \sum_{n\ge 2}\mathbb P(Y^n=\emptyset)\\
&=  \sum_{n\ge 2}\mathbb E[\mathbb P(Y^n=\emptyset|X^{n-1})]. 
\end{align*}
We note that in order for the process to die out in generation $n$ what is important is the {\em free degrees} of the points in $Y^{n-1}$ that is, for $Y_i\in Y^{n-1}$ we are interested in 
$$
N(Y_i, X^{n-1})=\#\{x\in \mathbb Z^2\colon x\in \partial Y_i\setminus X^{n-1}\},
$$
where $\partial Y_i$ denotes the set of $\mathbb Z^d$-neighbors of $Y_i$. We note that almost surely
\begin{align*}
\mathbb P(Y^n=\emptyset|X^{n-1})&=1\{N(Y_i, X^{n-1})\le 2\text{ for all }Y_i\in Y^{n-1}\}\mathbb P(Y^n=\emptyset|X^{n-1})
\end{align*}
and 
\begin{align*}
\mathbb P(Y^n=\emptyset|X^{n-1})= \Big(\frac{1}{6}\Big)^{\sum_{Y_i\in Y^{n-1}}1\{N(Y_i,X^{n-1})=2\}} \Big(\frac{1}{2}\Big)^{\sum_{Y_i\in Y^{n-1}}1\{N(Y_i,X^{n-1})=1\}}.
\end{align*}
Hence, the evolution of the probability of $Y^n=\emptyset$ is determined by the evolution of 
$$
\mu_{n-1}=\sum_{Y_i\in Y^{n-1}}\delta_{N(Y_i,X^{n-1})}\in \mathcal M(\{0,1,2,3\}),
$$
hence 
\begin{align*}
\mathbb P(Y^n=\emptyset|X^{n-1})&=1^{\mu_{n-1}(0)}\Big(\frac{1}{2}\Big)^{\mu_{n-1}(1)}\Big(\frac{1}{6}\Big)^{\mu_{n-1}(2)}0^{\mu_{n-1}(3)}\\
&=:\exp\Big(\sum_{k=1}^3\mu_{n-1}(k)g(k)\Big)=\exp(\mu_{n-1}(g)) =:f(\mu_{n-1})
\end{align*}
and $f$ is decreasing, i.e., for $\mu\le \nu$ (that is $\mu(k)\le \nu(k)$ for all $k\in \{0,\dots, 3\}$) we have $f(\mu)\ge f(\nu)$. In particular, 
\begin{align*}
1-\theta(2,2)&\le  \sum_{n\ge 1}\mathbb E[{\rm e}^{\mu_n(g)}],
\end{align*}
which should be bounded away from one. 
\color{blue} I was hoping for some kind of martingale argument but did not find one yet. -- Bene\color{black}

\begin{comment}{
\subsection{Undirected variants}
Instead of considering the oriented $k$-nearest neighbor graph itself, we can transform it into an undirected graph in an undirected fashion, i.e., connecting two vertices by an edge whenever there is a directed edge between them in at least one direction, or in a bidirectional fashion, i.e., connecting them if and only if there are directed edges between them in both directions. The following assertion is easy to verify. Let us call the former graph the $\mathbf U$-$k$nnG  and the latter one the $\mathbf B$-$k$nnG, respectively. It is trivial that the $\mathbf B$-$1$nnG does not percolate in any dimensions because its degrees are bounded by 1. Further, for $d=1$ and $k=2$, the graph percolates because it includes all nearest-neighbour edges of $\mathbb Z$.
%
\begin{prop} \color{red} obsolete \color{black}
The $\mathbf B$-2nnG does not percolate for $d\geq 2$.
\end{prop}
\begin{proof}
It suffices to verify the assertion for $d\geq 2$.
Since the $\mathbf B$-$k$nnG has degrees bounded by $k$, it follows that for $k=2$ any infinite path in the graph must be self-avoiding. Now, the probability that a fixed nearest-neighbor edge of $\mathbb Z^d$ is included in the $\mathbf B$-2nnG equals 
\[ \Big(\frac{1}{\binom{2d}{2}}\Big)^2 = \frac{1}{d^2}. \]
%
The number of self-avoided paths of length $n \in \mathbb N$ in $\Z^d$ starting from the origin is bounded by $(2d-1)^n$ (\color{red} cf.\ Grimmett \color{black}). 
Now, for a fixed self-avoiding path of length $2n$ consisting of nearest-neighbor edges of $\Z^d$ starting from $o$, consider the event that  
\begin{itemize}
\item for any $k=1,\ldots,n$, the $2k$-th edge of the path (i.e., the one between the $2k-1$-st and the $2k$-th vertex of the path) is included in the $\mathbf B$-2nnG (this has probability $1/\binom{2d}{2}=\frac{1}{d(2d-1)}$), and
\item for any $k=0,\ldots,n-1$, there is a \emph{directed} edge included in the 2-nnG from the $2k$-th vertex of the path to the $2k+1$-st one (this has probability $(2d-1)/\binom{2d}{2} = \frac{1}{d}$).  
\end{itemize}
Note that since the choices of outgoing 2nnG edges of different vertices are independent, the probability that there exists such a self-avoiding path is at most
\[ \frac{(2d-1)^{2n}}{d^{n} (d(2d-1))^n} = \big(\frac{2d-1}{d^2}\big)^n. \] 
Since $(2d-1)/d^2<1$ for any dimension $d\geq 2$, the proposition follows.
\end{proof}
%
%\bigskip
%Trying to repair our computation from last time, but now let $m$ be the number of corners in a blocking interface of length $l$, then we have $m-4$ anticorners and thus $l-2m+4$ vertices in the interface that are neither corners nor anticorners. Then we have
%\begin{align*}
%\mathbb P(\exists \text{ blocking interface } ) &\leq \sum_{n \geq 2} \mathbb P(\exists \text{ blocking interface through } (n,0) ) \\ &  \leq \sum_{n \geq 2} \sum_{l \geq 2n+2} \sum_{m \geq 1} \mathbb P(\exists \text{ blocking interface through } (n,0) \text{ with} \\ \quad &\text{$m$ corners and $l-2m+4$ vertices at straight lines}) 
%\end{align*}
%\color{red} Sorry, this doesn't seem to give a correct upper bound because requirements on blocking edges are often superfluous, as observed by Bas. -- András \color{black}
%
\section{Proofs}\label{sec:proof}
}
\end{comment}


\newpage 
\section{Branching process attempt \#2}
\textbf{Goal:} Use a branching process as upper-bound to conclude that we actually do not percolate in $d=k=2$. 

\subsection{Notation}
For $n \in \N$ denote by $S_n$ the $\ell^1$-ball in $\Z^2$ and by $\mathcal{R}_n$ the number of directed open paths in $S_n$ that connect the origin to the boundary $\partial S_n$.

\subsection{Heuristics}
If we have a directed open path $\gamma$ in $S_n$ that connects the origin to $\partial S_n$, then usually, its endpoint lies on one of the sides of $\partial S_n$ (and not in one of the four corners). Therefore, its offspring distribution $\mu$ is as follows: 
\begin{align*}
    \mu(0) = \frac{1}{6}, \quad \mu(1) = \frac{2}{3}, \quad \mu(2) = \frac{1}{6}.
\end{align*}
In particular, the typical offspring distribution is critical, and a Branching process with only these typical individuals would die out almost surely. Since it is exactly critical, it will survive for a substantial time which could explain that we don't see this in the simulations. 

Now if we have an open path in $S_n$ that reaches one of the corners of $S_n$, then it will have offspring distribution $\hat{\mu}$ given by 
\begin{align*}
    \mu(0) = 0, \quad \mu(1) = \frac{1}{2}, \quad \mu(2) = \frac{1}{2},
\end{align*}
which is super-critical. 

Now we need to make sure that this branching process dies out almost surely. 







\end{document}

\section{Continuum $k$-nearest-neighbor graphs}
\subsection{Expected degree}
\color{red} What I'm writing here is simple, it may or may not be useful in the end, but I want to write it down before I lose or forget it. -- Bene \color{black}

Let us calculate the expected degree of the typical point in the bidirectional nn-graph based on a Poisson point process
\begin{align*}
\mathbb E\Big[\sum_{X_i\in X}&1\{X(B_{|X_i|}(o))\le k\text{ and }X(B_{|X_i|}(X_i))\le k \}\Big]\\
&=\lambda \int_{\mathbb R^2}{\rm d} x \mathbb P(X(B_{|x|}(o))\le k\text{ and }X(B_{|x|}(x))\le k)\\
&=\lambda \int_0^\infty{\rm d} r 2\pi r \mathbb P(X(B_{r}(o))\le k\text{ and }X(B_{r}(r e_1))\le k)
\end{align*}
Let us note that replacing {\em and} by {\em or} we arrive at a formula for the undirected case. Further, consider the three areas $A_1=B_{r}(o)\setminus B_{r}(re_1)$,  $A_2=B_{r}(r e_1)\setminus B_{r}(o)$ and $B=B_{r}(o)\cap B_{r}(re_1)$ and note that 
$$|B|=r ^2(\frac{\pi}{3}-\frac{\sqrt{3}}{4})=:r^2c$$
and hence $|A_1|=|A_2|=r^2(\pi-c)$. Now let $V_1,V_2$ be two independent Poisson random variables with parameter $\lambda|A_i|$ and $W$ another independent Poisson random variable with parameter $\lambda|B|$. Then, 
\begin{align*}
\mathbb P(X(B_{r}(o))\le k\text{ and }X(B_{r}(r e_1))\le k)=\mathbb P(V_1+W\le k-1\text{ and }V_2+W\le k-1)
\end{align*}
Let us for example calculate the case $k=2$. 
\begin{align*}
\mathbb P(V_1+W\le 1\text{ and }V_2+W\le 1)=\sum_{n=0}^1 \mathbb P(W=n)\mathbb  P(V_1\le 1-n)^2\approx 1.5
\end{align*}

\subsection{Subcriticality for bidirectional nearest-neighbor graphs in high dimensions}
Heuristically, it should be difficult for any pair of points to choose each other as neighbors if the dimension is high. So, let us try a Gouere-type argument. Consider the event $G_\alpha$ that the origin is connected to $B_{10\alpha}\setminus B_{9\alpha}$ only using points in $B_{10\alpha}$. Later it will become important to have that $\mathbb P(G_\alpha)\le 1/2$ for $\alpha\in [1,10]$. Using that 
$$\mathbb P(G_{10\alpha})\le \mathbb P(G_{\alpha})\le 1-\exp(-\lambda|B_{\alpha}|)$$
and the fact that the percolation behavior of our model is independent of $\lambda$, we can choose $\lambda=\log 2/|B_\alpha|$ for a suitable $\alpha$ and this is satisfied.

In order for the argument to work, we need to make the probabilities of following error events small, uniformly in $\alpha$:
\begin{align*}
H_\alpha&=\{\text{ there exists a directly connected pair $x,y\in B_{100\alpha}$ with $|x-y|>\alpha$}\}\\
F_\alpha&=\{\text{ there exists a directly connected pair $x,y$ with $x\in B^c_{10\alpha}$ and $y\in B_{9\alpha}$}\}. 
\end{align*}
Using a first moment bound, we have that 
\begin{align*}
\mathbb P(H_\alpha)&\le \lambda^2\int_{B_{100\alpha}}d x\int_{B_{100\alpha}}d y1\{|x-y|>\alpha\}\mathbb P(x\leftrightarrow y)\\
\mathbb P(F_\alpha)&\le \lambda^2\int_{B^c_{10\alpha}}d x\int_{B_{9\alpha}}d y\mathbb P(x\leftrightarrow y)
\end{align*}
Note that $\mathbb P(x\leftrightarrow y)=f_k(\lambda|B_{|x-y|}|)$ since we are dealing with Poisson probabilities. Using the form of $\lambda$ we hence have that 
\begin{align*}
\mathbb P(H_\alpha)&\le c_1\int_{B_{100}}d x\int_{B_{100}}d y1\{|x-y|>1\}f_k(|B_{|x-y|}|)\\
\mathbb P(F_\alpha)&\le c_2\int_{B^c_{10}}d x\int_{B_{9}}d yf_k(|B_{|x-y|}|)
\end{align*}

\newpage

\section{Some further brainstorming}
From Bene:
\begin{itemize}
\item The B-kNN graph is related to SIR graphs by the fact that in both cases the edge-drawing mechanism depends not only on the relation of two points but also on other points
\item The B-kNN graph based on Poisson points should be stabilizing
\item One could consider other Poisson--Boolean models in which the ball around each point is random and depends on the surrounding point configuration.
\item The B-kNN graph is also somehow related to Poisson (directed) navigations.
\end{itemize}
Some more (from András):
\begin{itemize}
\item We see that even if we knew that there was an infinite cluster containing infinitely many pairwise disjoint infinite paths (or: ``the infinite cluster'' if we also had uniqueness), there must not be an essential trifurcation point, due to possible spider webs. But at least, there exists a degree-3 point with 3 edge-disjoint paths going out from it to infinity, call such a point \emph{good}. So there must be at least one infinite cluster with infinitely many good points. Now, I hope that for each good point we can find a spanning tree of the infinite cluster within which the good point is an essential trifurcation point. If we can even show that in any (nice) spanning tree the essential trifurcation points have a positive density, we can maybe get back to the Burton--Keane style argument that implies that the boundary of the graph in a ball/box is too large. \\
That is, I suspect that the spider-web is just some kind of decoration of an infinite 3-uniform tree.
\item In any case, such a line of argument requires the existence of an infinite cluster with infinitely many pairwise disjoint infinite paths, which I don't know how to prove, but I still hope that looking at half-spaces/sectors is the easiest way to attack this problem.
\end{itemize}

\subsection{A possible line of attack}
Let us try to reprove absence of percolation in the Bk-NN graph for $k=2$ via the following branching method. Let us denote by $B_r(x)$ the open ball with radius $r$ centered at $x$ and by $D_r(x)=B_r(x)\setminus \{x\}$, the ball with removed center. Then, we have for any $\alpha> 0$ we have
\begin{align*}
&\mathbb P(o\leftrightsquigarrow\infty)\le \mathbb P(\text{there exists a self-avoiding path starting at $o$ and ending in $B^{\rm c}_\alpha(o)$})\\
&\le \sum_{n\ge 1}\mathbb E\Big[\sum_{X_1,\dots, X_n\in X}^{\neq}\prod^n_{i=1}1\{X(D_{|X_i-X_{i-1}|}(X_{i-1}))\le 1,X(D_{|X_i-X_{i-1}|}(X_{i}))\le 1\}1\{X_n\in B^{\rm c}_\alpha(o)\}\Big]
\end{align*}
where $X_0\equiv o$ denotes the origin. The bound $\le 1$ indicates that in the open balls with removed centers, there can be at most one additional point. When moving to larger $k$, such as $k=3$, this bound changes accordingly. The idea is that for small $n$, the expectations are small because for large $\alpha$ we require large void spaces. On the other hand, for fixed $\alpha$ and large $n$, in order not to produce loops, the forbidden areas become also large. 

We can use the Slyvniak--Mecke formula to express, for any $n\ge 1$
\begin{align*}
&\mathbb E\Big[\sum_{X_1,\dots, X_n\in X}^{\neq}\prod^n_{i=1}1\{X(D_{|X_i-X_{i-1}|}(X_{i-1}))\le 1,X(D_{|X_i-X_{i-1}|}(X_{i}))\le 1\}1\{X_n\in B^{\rm c}_\alpha(o)\}\Big]\\
&=\int{\rm d} x_1\dots\int{\rm d} x_{n-1} \int_{B^{\rm c}_\alpha(o)}{\rm d} x_n\\
&\hspace{1cm}\mathbb E\Big[\prod^n_{i=1}1\{(X\cup\{x_1,\dots,x_n\})(D_{|x_i-x_{i-1}|}(x_{i-1}))\le 1,(X\cup\{x_1,\dots,x_n\})(D_{|x_i-x_{i-1}|}(x_{i}))\le 1\}\Big].
\end{align*}
Let us note that in the events $\{(X\cup\{x_1,\dots,x_n\})(D_{|x_i-x_{i-1}|}(x_{i-1}))\le 1\}$ on the right-hand side, the additional points $x_{i-1},x_i$ are irrelevant. This method, if successful, can also be used to say something about the decay of the finite-volume percolation function. 

Now it becomes tedious, 
\begin{align*}
&\mathbb E\Big[\prod^n_{i=1}1\{(X\cup\{x_1,\dots,x_n\})(D_{|x_i-x_{i-1}|}(x_{i-1}))\le 1,(X\cup\{x_1,\dots,x_n\})(D_{|x_i-x_{i-1}|}(x_{i}))\le 1\}\Big]\\
&=\mathbb E\Big[\Big(\sum_{A_1\subset [n]}\prod_{i\in A_1}1\{(\{x_1,\dots,x_n\})(D_{|x_i-x_{i-1}|}(x_{i-1}))=1\}1\{X(D_{|x_i-x_{i-1}|}(x_{i-1}))=0\}\times\\
&\qquad\qquad\qquad\prod_{i\in A^c_1}1\{(\{x_1,\dots,x_n\})(D_{|x_i-x_{i-1}|}(x_{i-1}))=0\}1\{X(D_{|x_i-x_{i-1}|}(x_{i-1}))\le1\}\Big)\times\\
&\qquad\Big(\sum_{A_2\subset [n]}\prod_{i\in A_2}1\{(\{x_1,\dots,x_n\})(D_{|x_i-x_{i-1}|}(x_{i}))=1\}1\{X(D_{|x_i-x_{i-1}|}(x_{i}))=0\}\times\\
&\qquad\qquad\qquad\prod_{i\in A^c_2}1\{(\{x_1,\dots,x_n\})(D_{|x_i-x_{i-1}|}(x_{i}))=0\}1\{X(D_{|x_i-x_{i-1}|}(x_{i}))\le1\}\Big)\Big]
\end{align*}
which is equal to 
\begin{align*}
&\sum_{A_1,A_2\subset [n]}\prod_{i\in A_1}1\{(\{x_1,\dots,x_n\})(D_{|x_i-x_{i-1}|}(x_{i-1}))=1\}\prod_{i\in A^c_1}1\{(\{x_1,\dots,x_n\})(D_{|x_i-x_{i-1}|}(x_{i-1}))=0\}\times\\
&\qquad\qquad\prod_{i\in A_2}1\{(\{x_1,\dots,x_n\})(D_{|x_i-x_{i-1}|}(x_{i}))=1\}\prod_{i\in A^c_2}1\{(\{x_1,\dots,x_n\})(D_{|x_i-x_{i-1}|}(x_{i}))=0\}\times\\
&\qquad\mathbb E\Big(\prod_{i\in A_1}1\{X(D_{|x_i-x_{i-1}|}(x_{i-1}))=0\}\prod_{i\in A^c_1}1\{X(D_{|x_i-x_{i-1}|}(x_{i-1}))\le1\}\times\\
&\qquad\qquad\prod_{i\in A_2}1\{X(D_{|x_i-x_{i-1}|}(x_{i}))=0\}\prod_{i\in A^c_2}1\{X(D_{|x_i-x_{i-1}|}(x_{i}))\le1\}\Big)
\end{align*}
So what happens here is that we prescribe deterministic points $x_1,\dots,x_n$ and via $A_1,A_2$ specify that for example $x_i$ is the second nearest neighbor of $x_{i-1}$ and $x_{i-1}$ is the first or second nearest neighbor of $x_{i}$. This sequence of deterministic points has to be legitimate in the sense that the points are in a configuration that is not impossible to achieve by putting additional (random) points. Now, given this sequence of prescribed points, we have to evaluate the probability to see the random points to be positioned precisely such that the sequence is legitimate. Let us note that that points may need to be places in order to make for example some deterministic point the second nearest neighbor. In the end, the probabilities can be evaluated by sums of independent Poisson random variables with parameters given by areas of intersecting balls. 

Next step:
\begin{align*}
&\mathbb E\Big(\prod_{i\in A_1}1\{X(D_{|x_i-x_{i-1}|}(x_{i-1}))=0\}\prod_{i\in A^c_1}1\{X(D_{|x_i-x_{i-1}|}(x_{i-1}))\le1\}\times\\
&\qquad\qquad\prod_{i\in A_2}1\{X(D_{|x_i-x_{i-1}|}(x_{i}))=0\}\prod_{i\in A^c_2}1\{X(D_{|x_i-x_{i-1}|}(x_{i}))\le1\}\Big)\\
&=\sum_{B_1\subset A_1^c}\sum_{B_2\subset A_2^c}\mathbb E\Big(\prod_{i\in A_1\cup B_1}1\{X(D_{|x_i-x_{i-1}|}(x_{i-1}))=0\}\prod_{i\in A^c_1\setminus B_1}1\{X(D_{|x_i-x_{i-1}|}(x_{i-1}))=1\}\times\\
&\qquad\qquad\prod_{i\in A_2\cup B_2}1\{X(D_{|x_i-x_{i-1}|}(x_{i-1}))=0\}\prod_{i\in A^c_2\setminus B_2}1\{X(D_{|x_i-x_{i-1}|}(x_{i-1}))=1\}\Big)\\
&=\sum_{B_1\subset A_1^c}\sum_{B_2\subset A_2^c}\mathbb E\Big(1\{X(\bigcup_{i\in A_1\cup B_1\cup A_2\cup B_2}D_{|x_i-x_{i-1}|}(x_{i-1}))=0\}\times\\
&\qquad\qquad \prod_{i\in (A^c_1\setminus B_1)\cup (A^c_2\setminus B_2)}1\{X(D_{|x_i-x_{i-1}|}(x_{i-1}))=1\}\Big)\\
&=\sum_{B_1\subset A_1^c}\sum_{B_2\subset A_2^c}\exp\Big(-|\bigcup_{i\in A_1\cup B_1\cup A_2\cup B_2}D_{|x_i-x_{i-1}|}(x_{i-1})|\Big)\times\\
&\qquad\qquad \mathbb E\Big[\prod_{i\in (A^c_1\setminus B_1)\cup (A^c_2\setminus B_2)}1\{X(D_{|x_i-x_{i-1}|}(x_{i-1})\setminus \bigcup_{j\in A_1\cup B_1\cup A_2\cup B_2}D_{|x_j-x_{j-1}|}(x_{j-1}))=1\}\Big]\\
\end{align*}
{\color{blue} Bas: Isn't the second line of the second equation incorrect? I think it should be $D_{|x_i-x_{i-1}|}(x_i)$ instead of $D_{|x_i-x_{i-1}|}(x_{i-1})$. And if so, the third and fourth step need correcting, too, I'd say.
}

\color{red}
A bit tighter upper bound, which might be a bit easier to compute because in many areas we just have void spaces after applying the Slivnyak--Mecke formula and thus we don't need much inclusion-exclusion (let me abuse notation and use $X$ both as a set and as the number of Poisson points in the set):
\begin{align*}
&\mathbb P(o\leftrightsquigarrow\infty)\le \mathbb P(\text{there exists a self-avoiding path starting at $o$ and ending in $B^{\rm c}_\alpha(o)$} \\ &\qquad \text{whose vertices do not form a loop in the whole $\mathbf B$-2NN graph} )\\
&\le \sum_{n\ge 1}\mathbb E\Big[\sum_{X_1,\dots, X_n\in X}^{\neq}\prod^n_{i=1}1\{X \cap D_{|X_i-X_{i-1}|}(X_{i-1}) \subseteq \{ X_{i-2} \}, \forall i=2,\ldots,n, \\ & X\cap D_{|X_i-X_{i-1}|}(X_{i}) \subseteq \{ X_{i+1} \}, \forall i=1,\ldots,n-1, \\& X(D_{|X_1-X_{0}|}(X_0)) \leq 1, X_n \notin D_{|X_1-X_{0}|}(X_i) , X(D_{|X_n-X_{n-1}|}(X_n)) \leq 1, X_0 \notin D_{|X_n-X_{n-1}|}(X_n) \} \Big]
\end{align*}
Using the Slivnyak--Mecke formula, for fixed $n$ we can express the expectation as
\begin{align*}
\int& {\rm d} x_1 \ldots \int {\rm d} x_{n-1} \int_{B_\alpha^c(o)} {\rm d} x _n \\ & \mathds E \Big[ \prod_{i=2}^n \mathds 1 \big\{ (X \cup \{ x_1,\ldots,x_n\}) (D_{|x_i-x_{i-1}|}(x_{i-1}) \subseteq \{ x_{i-2} \} \big\} \\
&\times \prod_{i=1}^{n-1} \mathds 1 \big\{ \{ (X \cup \{ x_1,\ldots,x_n\}) (D_{|x_i-x_{i-1}|}(x_{i}) \subseteq \{ x_{i+1} \} \big\} \\&
\qquad \times \mathds 1 \{ x_n \notin (X \cup \{ x_1,\ldots,x_n\}) (D_{|x_1-x_{0}|}(x_{0}), (X \cup \{ x_1,\ldots,x_n\}) (D_{|x_1-x_{0}|}(x_{0}) \leq 1 \} \\ &\qquad \times \mathds 1 \{ o \notin (X \cup \{ x_1,\ldots,x_n\}) (D_{|x_n-x_{n-1}|}(x_{n-1}), (X \cup \{ x_1,\ldots,x_n\}) (D_{|x_n-x_{n-1}|}(x_{n-1}) \leq 1 \} \Big].
\end{align*}

\newpage


\section{General graphs: trifurcation}
\color{red} What I'm writing here is simple, it may or may not be useful in the end, but I want to write it down before I lose or forget it. -- András \color{black}

Let $G=(V,E)$ be a countably infinite connected simple graph with degrees bounded by 3. An \emph{essential trifurcation point} is a vertex $v \in G$ such that the graph with vertex set $V \setminus \{ v \}$ and edge set $E|_{(V \setminus \{ v\})^2}$ consists of three infinite connected components. Of course, essential trifurcation points have degree 3.

\begin{lemma}\label{lem-loop}
No essential trifurcation point of $G$ is included in a loop of the graph.
\end{lemma}
\begin{proof}
Let $v \in V$ be a degree-3 point included in a loop. Let $v_1,v_2,v_3$ be its neighbours. At least two of these, say $v_1,v_2$, are included in the loop. Hence, there exists a path between $v_1$ and $v_2$ not including $v$. But then, after removing $v$ and the edges $(v,v_1),(v,v_2),(v,v_3)$ from the graph, $v_1$ and $v_2$ are still in the same connected component. Hence, $v$ cannot be an essential trifurcation point.
\end{proof}
\begin{cor}\label{cor-almosttree}
Let $v,w \in G$ be two different essential trifurcation points. Then there exists a unique pair $(v_1,w_1) \in G\times G$ such that any path between $v$ and $w$ starts with the edge $(v,v_1)$ and ends with the edge $(w_1,w)$.
\end{cor}
Note that the case when $v_1=w$ and $w_1=v$ is not excluded here.
\begin{proof}[Proof of Corollary~\ref{cor-almosttree}]
Assume otherwise, so that without loss of generality there exist $v_1,v_2$ adjacent to $v$ and one path from $v$ to $w$ starting with the edge $(v,v_1)$ and one starting with the edge $(v,v_2)$. If these two paths are edge-disjoint, then the union of their edges forms a loop including both $v$ and $w$, which contradicts with Lemma~\ref{lem-loop}. Otherwise, the two paths have a joint vertex $z$. Thus, the path consisting of the edges of the path from $v$ to $w$ starting with $(v,v_1)$ up to $z$ and the edges of the path from $z$ to $v$ form a loop including $v$, another contradiction with the lemma.
\end{proof}
Iterating this corollary, we obtain the following assertion.
\begin{cor}
Let $v,w$ be two different essential trifurcation points of $G$. Then, there exists a unique $n\in\mathbb N$ and $(v_0, v_1,\ldots,v_n) \in V^{n+1}$ mit $v_0=v$ and $w_0=w$ such that all the $v_i$ are essential trifurcation points and any path between $v$ and $w$ goes through the points $v_0,v_1,\ldots,v_n$ in this order.
\end{cor}

%Thus, for any two different essential trifurcation points $v,w \in V$ there exists a unique pair of vertices $(s(v,w),t(v,w)) \in G\times G$ such that any path from $v$ to $w$ starts with the edge $(v,s(v,w))$ and ends with the edge $(t(v,w),w)$. Of course, $s(v,w)=t(w,v)$ and $s(w,v)=t(v,w)$. We say that $v,w$ are \emph{almost direct neighbours} if there exists a path from $v$ to $w$ that does not contain any essential trifurcation point apart from $v$ and $w$. 

The \emph{essential tree} $T(G)=(V(T(G)),E(T(G)))$ of $G$ is the graph with vertex set $V(T(G))=\{v \in V \colon v \text{ essential trifurcation point of $G$}\}$ and edge set $E(T(G))=\{ (v,w) \in V(T(G)) \times V(T(G)) \colon v,w\text{ direct neighbours}\}$. 
\begin{lemma}\label{lem-contractedtree}
$T(G)$ is a tree.
\end{lemma}
\begin{proof}
The connectedness of $T(G)$ follows from the fact that for any $v,w \in V(T(G))$, the corresponding path $(v_0, v_1,\ldots,v_n)$ in the statement of Corollary~\ref{cor-almosttree} exists, and the loop-freeness of $T(G)$ follows from the uniqueness of the fact that for any $v,w \in V(T(G))$, the corresponding path is unique.
\end{proof}
\begin{lemma}\label{lem-tree}
Let $H=(V(H),E(H))$ be a tree with degrees bounded by 3, and let $F \subset V(H)$ be a collection of vertices of degree equal to 3. Now, if $|F|=k \in \mathbb N$, then $H$ has at least $2k+2$ vertices. Further, there are at least $\lceil (k+2)/2 \rceil$ vertices in $F$ that are connected by an edge to at least one point in $V(H) \setminus F$.
\end{lemma}
\begin{proof}
Let us call $v,w \in F$ direct neighbours if the (unique) path between them contains no degree-3 point apart from $v$ and $w$. Then, contracting paths between direct neighbours to a single edge, we cannot increase the number of vertices. After all these contractions, the edges between points of $F$ form a tree, which thus has $k-1$ vertices, and each vertex of $F$ still has degree 3. Now, there are $3k$ half-edges adjacent to vertices in $F$, and at most $2k-2$ of them can correspond to edges between vertices of $F$. Hence, there must still be at least $k+2$ half-edges that are adjacent to vertices outside $F$ that have not been contracted. These $k+2$ vertices must be pairwise distinct. Indeed, otherwise some of them , say $z$, would be adjacent to two vertices $v,w$ of $F$, and then $(v,z),(z,w)$ would form a path between $v$ and $z$ in $H$.
This is however impossible because we have another path between $v$ and $w$ in $H$ that only goes through points of $F$ and contracted vertices. We conclude that $H$ must have at least $|F|+k+2=2k+2$ vertices, as wanted.

The same contraction procedure can also not increase the number of vertices in $F$ adjacent to a point in $V(H) \setminus F$. After the contractions, there are precisely $k+2$ vertices outside $F$ adjacent to $F$, as we have seen, while each vertex $v$ of $F$ is connected by at least one edge to $F \setminus \{ v \}$, consequently it is connected to at most 2 vertices outside $F$. By the pigeon-hole principle, we conclude that at least $\lceil (k+2)/2 \rceil$ vertices of $F$ have to be adjacent to vertices outside $F$ in the contracted graph, and thus also in $H$.
\end{proof}
Let us now get back to our infinite graph $G$ studied before. For $W \subseteq V(G)$ define the \emph{inner boundary} of $W$ as
\[ \partial^{\text{in}}~W = \{ w \in W \colon \exists v \in V(G) \setminus W \colon (w,v) \in E(G) \} \]
and
\[ \partial^{\text{out}}~ W = \{ w \in V(G) \setminus W \colon \exists V\in W \colon (w,v) \in E(G) \}. \]
\begin{cor}
Let $W \subset V(G)$ be a set containing only essential trifurcation points of $G$ with $|V| \geq k$. Then, we have
\begin{enumerate}
\item $|\partial^{\text{in}}~W| \geq \lceil (k+2)/2 \rceil$,
\item $|\partial^{\text{out}}~ W| \geq k+2$.
\end{enumerate}
\end{cor}
\begin{proof}
The corollary is a direct consequence of Lemmas~\ref{lem-contractedtree} and \ref{lem-tree}. The tree $F$ corresponding to Lemma~\ref{lem-tree} can be chosen as any finite subtree of $T(G)$ whose vertex set contains $W$.
\end{proof}
\color{red} This should eventually yield us a ``linearly growing boundary'', although I'm not sure yet what to do with those points that are not essential trifurcation points. \color{black}


%\newpage
%\color{black}
%\subsection{Trying to repair the trifurcation argument}
%
%Back to the B-3NN graph, assume that there is an infinite cluster containing infinitely many edge-disjoint infinite paths (this would be true if we knew that infinite clusters ``grow in all directions''). Now, it is easy to show that this cluster contains a \emph{good} point, i.e., a point that has 3 edge-disjoint \color{blue} infinite \color{black} paths coming out of it. Consequently, good points must have a positive density (on the event of percolation if it is not a 0-1 event etc.). Now, we inductively grow a sequence of trees $(T_n)_{n\in\mathbb N_0}$ with infinitely many essential trifurcation points that is a subgraph of the original infinite cluster.
%
%Consider a good point $w_0$ and the three edge-disjoint infinite paths $p_1,p_2,p_3$ going out of it. Initially, $T_0$ consists of the vertices and edges of the union of these three paths. Since the infinite cluster is connected and contains infinitely many edge-disjoint infinite paths, there must exist a good point $w$ in one of the three paths (apart from $w_0$), say in $p_1$. Two edges adjacent to $w$ belong to $p_1$, let $e$ denote the third edge. The question is now whether there exists an infinite path adjacent to $w$ that is edge-disjoint from $p_2$ and $p_3$. Note that if there exists one, it also has to be vertex-disjoint from $p_2$ and $p_3$ because of the degree bound of 3 (and because $w$ is not equal to $w_0$). So if there exists none, this is because all paths starting from $w$ with $e$ share at least one edge with $p_2$ or $p_3$. Otherwise, from $w$ there is a path to infinity starting with $e$ that is edge-disjoint from $p_1,p_2,p_3$. 
%
%Since the cluster contains infinitely many edge-disjoint infinite paths, we conclude that on $p_1,p_2$ or $p_3$ there exists a good point from which there is a path to infinity that is edge-disjoint from $p_1,p_2,p_3$; choose such a point $w_1$ and add the edges and remaining vertices of the infinite path, say $p_4$, to $T_0$. The graph $T_1$ we obtain this way is still a tree.
%
%Iterating this procedure, after $n$ steps we can always find a new point $w_n$ in the current tree $T_n$ from which an infinite path goes out to $\infty$, being edge-disjoint from all the previous $n+2$ infinite paths. After $n$ steps, we have $n$ edge-disjoint infinite paths in $T_n$, and $T_n$ is a tree. So we can take the union $T$ of all these trees. By construction we obtain a tree \color{red} clear? Does connectedness need a proof? Acyclicity needs no proof \color{black}, and the points $w_n$ are always essential trifurcation points of this tree because if one considers two points $x,y$ on two different edge-disjoint paths starting from $w_n$, then for $k\geq n$ the only path in $T_k$ from $x$ to $y$ is via $w_n$, so that if we remove $w_n$, $T_k$ splits into three infinite clusters, and thus so does $T$. 
%
%\color{red} Now, if this really work, we just have to show that conditional on the existence of an infinite cluster, there exists an infinite cluster containing infinitely many infinite paths. The rest should follow from the Burton--Keane argument... \color{black}
\end{document}



%%%%%%%%%%%%%%%%% References %%%%%%%%%%%%%%%%%%
\footnotesize{\printbibliography}
\end{document}
